\documentclass[10pt, landscape, a4paper, onesided]{extarticle}
\usepackage{xcolor}

\definecolor{pupil}{RGB}{0, 128, 10}
\definecolor{specialist}{RGB}{3, 168, 158}
\definecolor{expert}{RGB}{3, 168, 158}
\newcommand{\universityName}{Daffodil International University}
\newcommand{\teamName}{DIU\_YinYang.exe}
\newcommand{\firstMember}{\color{specialist}AnupBarman}
\newcommand{\secondMember}{\color{specialist}SiyamBhuiyan}
\newcommand{\thirdMember}{\color{specialist}NoObMin}

\usepackage{minted}
\usepackage{sectsty}
\usepackage{fouriernc}
\usepackage[T1]{fontenc}
\usepackage{framed}
\usepackage{makeidx}
\usepackage{enumitem}
\usepackage{titlesec}
\usepackage[scaled=0.9]{DejaVu Sans Mono}
\usepackage[utf8]{inputenc}
\usepackage[english]{babel}
\usepackage{listings}
\usepackage{amsmath}
\usepackage{verbatim}
\usepackage{hyperref}
\usepackage[dvipsnames]{xcolor}
\usepackage{multicol}
\usepackage{geometry}
\usepackage{amssymb,tikz,psfrag,epstopdf}
\geometry{verbose,landscape,a4paper,tmargin=1.25cm,bmargin=0cm,lmargin=0.75cm,rmargin=.25cm}

% Fix for overlapping subsection numbers in ToC
\usepackage{tocloft}
\setlength{\cftsubsecnumwidth}{3em}

\newcommand{\maketeampage}{
	\pagestyle{empty}
	\setcounter{page}{0}
	\begin{center}
		\strut % Magic
		\vspace{-1cm}\\
		\includegraphics[width=12cm]{logo}\\
		\vspace{-2cm}
		{\fontsize{24}{30} \fontfamily{ptm} \selectfont {\universityName}\\}
		\vspace{2cm}
		{\fontsize{60}{60} \fontfamily{ptm} \selectfont {\teamName}\\}
		\vspace{1cm}
		{\fontsize{24}{30} \fontfamily{ptm} \selectfont { {\firstMember}, {\secondMember}, {\thirdMember} }\\}
		\vspace{2cm}
		{\fontsize{24}{30} \fontfamily{ptm} \selectfont {Team Reference Document}}\\
		\vspace{1cm}
		\vspace{1cm}
	\end{center}
	\clearpage
	\pagestyle{fancy}
}

\newcommand{\printdescription}[1]{
  \codeheader{Description}{#1}
}
\newcommand{\printusage}[1]{
  \codeheader{Usage}{\texttt{#1}}
}
\newcommand{\printtime}[1]{
  \codeheader{Time}{#1}
}
\newcommand{\printmemory}[1]{
  \codeheader{Memory}{#1}
}
\newcommand{\codeheader}[2]{
	{\par\noindent\normalsize{\textbf{#1:} #2\par}}
}
\newcommand{\bigo}[1] {\ensuremath{\mathcal{O} \left( #1 \right)}}

\newenvironment{myitemize}
{ \begin{itemize}[leftmargin=.5cm]
    \setlength{\itemsep}{0pt}
    \setlength{\parskip}{0pt}
    \setlength{\parsep}{0pt}     }
{ \end{itemize}                  }

\setlength{\parindent}{0pt}
\newcommand*{\Scale}[2][4]{\scalebox{#1}{$#2$}}%

\setlength{\columnsep}{1.5mm}
\setlength{\columnseprule}{1px}

\setcounter{secnumdepth}{4}
\setcounter{tocdepth}{2}

\BeforeBeginEnvironment{minted}{\vspace{0mm}}
\AfterEndEnvironment{minted}{\vspace{-1mm}}

\linespread{0.30}
\usepackage{fancyhdr}
\pagestyle{fancy}
\fancyhf{}
\renewcommand{\headrulewidth}{1pt}
\renewcommand{\theFancyVerbLine}{\sffamily \textcolor[rgb]{0.15,0.15,0.15}{\small \arabic{FancyVerbLine}}}

\lhead{\textcolor{MidnightBlue}{\textit{\textbf{\universityName}}} | \textcolor{MidnightBlue}{\textit{\textbf{\teamName}}}}
\setlength\headheight{16pt}
\rhead{\thepage}
\headsep = 0pt
\textheight = 7.95in
\footskip = 0mm

\titlespacing*{\section}{0pt}{0}{0}
\titlespacing*{\subsection}{0pt}{0}{0}
\titlespacing*{\subsubsection}{0pt}{0}{0}

\titleformat{\section}
  {\normalfont\fontsize{10}{15}\bfseries}{\thesection}{1em}{}
\titleformat{\subsection}
  {\normalfont\fontsize{8}{15}\bfseries}{\thesubsection}{1em}{}
\titleformat{\subsubsection}
  {\normalfont\fontsize{8}{15}\bfseries}{\thesubsubsection}{1em}{}
\subsectionfont{\Huge}

\begin{document}
\maketeampage
\begin{multicols*}{3}
\tableofcontents
\vspace{2mm}
\noindent\rule{9.7cm}{1pt}
\vspace{2mm}\section{Setup}
\subsection{$CP\_Ubuntu}
\begin{minted}[
    breaklines = true,
    breakanywhere = true,
    frame=lines,    
    fontfamily=tt,
    linenos=false,
    numberblanklines=true,
    numbersep=2pt,
    gobble=0,
    framerule=1pt,
    framesep=1mm,
    funcnamehighlighting=true,
    tabsize=2,
    obeytabs=false,
    mathescape=false
    samepage=false, %with this setting you can force the list to appear on the same page
    showspaces=false,
    showtabs =false,
    texcl=false
]{C++}

{
	"cmd": ["ulimit -s 268435456; g++ -std=c++20 $file_name -o $file_base_name && timeout 4s ./$file_base_name < inputf.in > outputf.in"],
	"selector": "source.cpp",
	"shell": true,
	"working_dir": "$file_path"
}
\end{minted}

\subsection{$CP\_Windows}
\begin{minted}[
    breaklines = true,
    breakanywhere = true,
    frame=lines,    
    fontfamily=tt,
    linenos=false,
    numberblanklines=true,
    numbersep=2pt,
    gobble=0,
    framerule=1pt,
    framesep=1mm,
    funcnamehighlighting=true,
    tabsize=2,
    obeytabs=false,
    mathescape=false
    samepage=false, %with this setting you can force the list to appear on the same page
    showspaces=false,
    showtabs =false,
    texcl=false
]{C++}

{
	"cmd": ["g++.exe","-std=c++20", "${file}", "-o", "${file_base_name}.exe", "&&" , "${file_base_name}.exe<inputf.in>outputf.in"],
	"selector": "source.cpp",
	"shell": true,
	"working_dir": "$file_path"
}
\end{minted}

\subsection{$Stress Testing(check.sh)}
\begin{minted}[
    breaklines = true,
    breakanywhere = true,
    frame=lines,    
    fontfamily=tt,
    linenos=false,
    numberblanklines=true,
    numbersep=2pt,
    gobble=0,
    framerule=1pt,
    framesep=1mm,
    funcnamehighlighting=true,
    tabsize=2,
    obeytabs=false,
    mathescape=false
    samepage=false, %with this setting you can force the list to appear on the same page
    showspaces=false,
    showtabs =false,
    texcl=false
]{C++}

// chmod u+x check.sh
// ./check.sh
set -e
g++ gen.cpp -o gen
g++ code.cpp -o code
g++ brute.cpp -o brute
for ((i = 1; ; ++i)); do
	echo "Passed on TestCase: " $i
	./gen $i > in
	./code < in > out1
	./brute < in > out2
	diff -Z out1 out2 || break
done
echo -e "WA on the following test:"
cat in
echo -e "\nExpected:"
cat out2
echo -e "\nFound:"
cat out1
\end{minted}

\subsection{$Stress Testing(gen.cpp)}
\begin{minted}[
    breaklines = true,
    breakanywhere = true,
    frame=lines,    
    fontfamily=tt,
    linenos=false,
    numberblanklines=true,
    numbersep=2pt,
    gobble=0,
    framerule=1pt,
    framesep=1mm,
    funcnamehighlighting=true,
    tabsize=2,
    obeytabs=false,
    mathescape=false
    samepage=false, %with this setting you can force the list to appear on the same page
    showspaces=false,
    showtabs =false,
    texcl=false
]{C++}

#include <bits/stdc++.h>
using namespace std;
using ll = long long;

mt19937_64 rng(chrono::steady_clock::now().time_since_epoch().count());

inline ll gen_random(ll l, ll r) {
  return uniform_int_distribution<ll>(l, r)(rng);
}
inline double gen_random_real(double l, double r) {
  return uniform_real_distribution<double>(l, r)(rng);
}
int main(int argc, char* args[]) {
  int _ = atoi(args[1]);
  rng.seed(_);
  int n = gen_random(1, 5);
  vector<int> per;
  for (int i = 0; i < n; ++i) {
    per.push_back(i + 1);
  }
  shuffle(per.begin(), per.end(), rng);
  return 0;
}
\end{minted}

\section{Data Structures}
\subsection{Fast Unordered Map}
\begin{minted}[
    breaklines = true,
    breakanywhere = true,
    frame=lines,    
    fontfamily=tt,
    linenos=false,
    numberblanklines=true,
    numbersep=2pt,
    gobble=0,
    framerule=1pt,
    framesep=1mm,
    funcnamehighlighting=true,
    tabsize=2,
    obeytabs=false,
    mathescape=false
    samepage=false, %with this setting you can force the list to appear on the same page
    showspaces=false,
    showtabs =false,
    texcl=false
]{C++}

mp.reserve(1<<20);       // about 1M buckets
mp.max_load_factor(0.7); // safe and fast
\end{minted}

\subsection{Mex of All Subarray}
\begin{minted}[
    breaklines = true,
    breakanywhere = true,
    frame=lines,    
    fontfamily=tt,
    linenos=false,
    numberblanklines=true,
    numbersep=2pt,
    gobble=0,
    framerule=1pt,
    framesep=1mm,
    funcnamehighlighting=true,
    tabsize=2,
    obeytabs=false,
    mathescape=false
    samepage=false, %with this setting you can force the list to appear on the same page
    showspaces=false,
    showtabs =false,
    texcl=false
]{C++}

const int N = 1e5 + 9, inf = 1e9;
struct ST {
  int t[4 * N];
  ST() {}
  void build(int n, int b, int e) {
    t[n] = 0;
    if (b == e) {
      return;
    }
    int mid = (b + e) >> 1, l = n << 1, r = l | 1;
    build(l, b, mid);
    build(r, mid + 1, e);
    t[n] = min(t[l], t[r]);
  }
  void upd(int n, int b, int e, int i, int x) {
    if (b > i || e < i) return;
    if (b == e && b == i) {
      t[n] = x;
      return;
    }
    int mid = (b + e) >> 1, l = n << 1, r = l | 1;
    upd(l, b, mid, i, x);
    upd(r, mid + 1, e, i, x);
    t[n] = min(t[l], t[r]);
  }
  int get_min(int n, int b, int e, int i, int j) {
    if (b > j || e < i) return inf;
    if (b >= i && e <= j) return t[n];
    int mid = (b + e) >> 1, l = n << 1, r = l | 1;
    int L = get_min(l, b, mid, i, j);
    int R = get_min(r, mid + 1, e, i, j);
    return min(L, R);
  }
  int get_mex(int n, int b, int e, int i) {  // mex of [i... cur_id] if (b == e) return b;
    int mid = (b + e) >> 1, l = n << 1, r = l | 1;
    if (t[l] >= i) return get_mex(r, mid + 1, e, i);
    return get_mex(l, b, mid, i);
  }
} t;
int a[N], f[N];
int32_t main() {
  ios_base::sync_with_stdio(0);
  cin.tie(0);
  int n;
  cin >> n;
  for (int i = 1; i <= n; i++) {
    cin >> a[i];
    --a[i];
  }
  t.build(1, 0, n);
  set<array<int, 3>> seg;  // for cur_id = i, [x[0]... i], [x[0] + 1...i], ...[x[1]... i] has mex,→ ,→ x[2]
  for (int i = 1; i <= n; i++) {
    int x = a[i];
    int r = min(i - 1, t.get_min(1, 0, n, 0, x - 1));
    int l = t.get_min(1, 0, n, 0, x) + 1;
    if (l <= r) {
      auto it = seg.lower_bound({l, -1, -1});
      while (it != seg.end() && (*it)[1] <= r) {
        auto x = *it;
        it = seg.erase(it);
      }
    }
    t.upd(1, 0, n, x, i);
    for (int j = r; j >= l;) {
      int m = t.get_mex(1, 0, n, j);
      int L = max(l, t.get_min(1, 0, n, 0, m) + 1);
      f[m] = 1;
      seg.insert({L, j, m});
      j = L - 1;
    }
    int m = !a[i];
    seg.insert({i, i, m});
    f[m] = 1;
  }
  int ans = 0;
  while (f[ans]) ++ans;
  cout << ans + 1 << '\n';
  return 0;
}
\end{minted}

\subsection{Pbds}
\begin{minted}[
    breaklines = true,
    breakanywhere = true,
    frame=lines,    
    fontfamily=tt,
    linenos=false,
    numberblanklines=true,
    numbersep=2pt,
    gobble=0,
    framerule=1pt,
    framesep=1mm,
    funcnamehighlighting=true,
    tabsize=2,
    obeytabs=false,
    mathescape=false
    samepage=false, %with this setting you can force the list to appear on the same page
    showspaces=false,
    showtabs =false,
    texcl=false
]{C++}

#include <ext/pb_ds/assoc_container.hpp>
#include <ext/pb_ds/tree_policy.hpp>
#include <functional>
using namespace __gnu_pbds;
typedef tree<int, null_type, less<int>, rb_tree_tag,
             tree_order_statistics_node_update>
    ordered_set;
// s.order_of_key(x) = number of elements strictly less than x
// *s.find_by_order(i) = ith element in set (0 index)
\end{minted}

\subsection{Segment Tree(BSUA)}
\begin{minted}[
    breaklines = true,
    breakanywhere = true,
    frame=lines,    
    fontfamily=tt,
    linenos=false,
    numberblanklines=true,
    numbersep=2pt,
    gobble=0,
    framerule=1pt,
    framesep=1mm,
    funcnamehighlighting=true,
    tabsize=2,
    obeytabs=false,
    mathescape=false
    samepage=false, %with this setting you can force the list to appear on the same page
    showspaces=false,
    showtabs =false,
    texcl=false
]{C++}

// CSES - 1749
const int MX = 2e5 + 10;
int n;
int arr[MX], st[MX << 2];
void assign(int i, int x, int u = 1, int s = 0, int e = n - 1) {
  if (s == e) {
    st[u] = x;
    return;
  }
  int v = u << 1, w = v | 1, m = (s + e) >> 1;
  if (i <= m) assign(i, x, v, s, m);
  else assign(i, x, w, m + 1, e);
  st[u] = st[v] + st[w];
}
int kth(int k, int u = 1, int s = 0, int e = n - 1) {
  if (st[u] < k) return -1;
  if (s == e) {
    return s;
  }
  int v = u << 1, w = v | 1, m = (s + e) >> 1;
  if (st[v] >= k) return kth(k, v, s, m);
  else return kth(k - st[v], w, m + 1, e);
}
void solve() {
  cin >> n;
  for (int i = 0; i < n; ++i) {
    cin >> arr[i];
  }
  for (int i = 0; i < n; ++i) {
    assign(i, 1);
  }
  for (int i = 0; i < n; ++i) {
    int x;
    cin >> x;
    int ind = kth(x);
    assign(ind, 0);
    cout << arr[ind] << " ";
  }
}
\end{minted}

\subsection{Segment Tree(LzP)}
\begin{minted}[
    breaklines = true,
    breakanywhere = true,
    frame=lines,    
    fontfamily=tt,
    linenos=false,
    numberblanklines=true,
    numbersep=2pt,
    gobble=0,
    framerule=1pt,
    framesep=1mm,
    funcnamehighlighting=true,
    tabsize=2,
    obeytabs=false,
    mathescape=false
    samepage=false, %with this setting you can force the list to appear on the same page
    showspaces=false,
    showtabs =false,
    texcl=false
]{C++}

class stree {
  vector<ll> st, lazy;

 public:
  stree(int n) {
    st.assign((n << 2) + 10, 0);
    lazy.assign((n << 2) + 10, 0);
  }
  void push(int u, int s, int e) {
    if (!lazy[u]) return;
    st[u] += (e - s + 1) * 1LL * lazy[u];
    if (s != e) {
      int v = u << 1, w = v | 1, m = (s + e) >> 1;
      lazy[v] += lazy[u];
      lazy[w] += lazy[u];
    }
    lazy[u] = 0;
  }
  void build(int u, int s, int e, int arr[]) {
    if (s == e) {
      st[u] = arr[s];
      return;
    }
    int v = u << 1, w = v | 1, m = (s + e) >> 1;
    build(v, s, m, arr);
    build(w, m + 1, e, arr);
    st[u] = st[v] + st[w];
  }
  void update(int l, int r, int x, int u, int s, int e) {
    push(u, s, e);
    if (e < l or r < s) return;
    int v = u << 1, w = v | 1, m = (s + e) >> 1;
    if (l <= s and e <= r) {
      st[u] += (e - s + 1) * 1LL * x;
      if (s != e) {
        lazy[v] += x;
        lazy[w] += x;
      }
      return;
    }
    update(l, r, x, v, s, m);
    update(l, r, x, w, m + 1, e);
    st[u] = st[v] + st[w];
  }
  ll query(int l, int r, int u, int s, int e) {
    push(u, s, e);
    if (e < l or r < s) return 0;
    if (l <= s and e <= r) return st[u];
    int v = u << 1, w = v | 1, m = (s + e) >> 1;
    return query(l, r, v, s, m) + query(l, r, w, m + 1, e);
  }
};
\end{minted}

\subsection{Segment Tree}
\begin{minted}[
    breaklines = true,
    breakanywhere = true,
    frame=lines,    
    fontfamily=tt,
    linenos=false,
    numberblanklines=true,
    numbersep=2pt,
    gobble=0,
    framerule=1pt,
    framesep=1mm,
    funcnamehighlighting=true,
    tabsize=2,
    obeytabs=false,
    mathescape=false
    samepage=false, %with this setting you can force the list to appear on the same page
    showspaces=false,
    showtabs =false,
    texcl=false
]{C++}

class stree {
  vector<ll> st;

 public:
  stree(int n) {
    st.assign((n << 2) + 10, 0);
  }
  void build(int u, int s, int e, int arr[]) {
    if (s == e) {
      st[u] = arr[s];
      return;
    }
    int v = u << 1, w = v | 1, m = (s + e) >> 1;
    build(v, s, m, arr);
    build(w, m + 1, e, arr);
    st[u] = st[v] + st[w];
  }
  ll query(int l, int r, int u, int s, int e) {
    if (e < l or r < s) return 0;
    if (l <= s and e <= r) return st[u];
    int v = u << 1, w = v | 1, m = (s + e) >> 1;
    return query(l, r, v, s, m) + query(l, r, w, m + 1, e);
  }
  void update(int i, int x, int u, int s, int e) {
    if (s == e) {
      st[u] = x;
      return;
    }
    int v = u << 1, w = v | 1, m = (s + e) >> 1;
    if (i <= m) update(i, x, v, s, m);
    else update(i, x, w, m + 1, e);
    st[u] = st[v] + st[w];
  }
};
\end{minted}

\subsection{Sparse Table}
\begin{minted}[
    breaklines = true,
    breakanywhere = true,
    frame=lines,    
    fontfamily=tt,
    linenos=false,
    numberblanklines=true,
    numbersep=2pt,
    gobble=0,
    framerule=1pt,
    framesep=1mm,
    funcnamehighlighting=true,
    tabsize=2,
    obeytabs=false,
    mathescape=false
    samepage=false, %with this setting you can force the list to appear on the same page
    showspaces=false,
    showtabs =false,
    texcl=false
]{C++}

const int mxN = 1e5 + 10, M = 21;
int sparse[mxN][M];
void build_sparse(int n, vector<int>& v) {
  for (int i = 0; i < n; i++) sparse[i][0] = v[i];
  for (int k = 1; k < M; k++) {
    for (int i = 0; i + (1 << k) <= n; i++) {
      sparse[i][k] = max(sparse[i][k - 1], sparse[i + (1 << (k - 1))][k - 1]);
    }
  }
}
int query(int l, int r) {  // 0 based index
  if (l > r) swap(l, r);
  int b = __bit_width(r - l + 1) - 1;
  return max(sparse[l][b], sparse[r - (1 << b) + 1][b]);
}
\end{minted}

\section{Dynamic Programming}
\subsection{CoinChange(MinCoins)}
\begin{minted}[
    breaklines = true,
    breakanywhere = true,
    frame=lines,    
    fontfamily=tt,
    linenos=false,
    numberblanklines=true,
    numbersep=2pt,
    gobble=0,
    framerule=1pt,
    framesep=1mm,
    funcnamehighlighting=true,
    tabsize=2,
    obeytabs=false,
    mathescape=false
    samepage=false, %with this setting you can force the list to appear on the same page
    showspaces=false,
    showtabs =false,
    texcl=false
]{C++}

const int INF = 1e9;
int min_coins(int target, vector<int>& coins) {
  // DP Definition: dp[i] = minimum coins needed to make value i
  vector<int> dp(target + 1, INF);
  // Base Case: 0 coins are needed to make value 0
  dp[0] = 0;
  for (int i = 1; i <= target; i++) {
    for (int c : coins) {
      if (i - c >= 0) {
        // Transition: Try taking coin 'c', adding 1 to the count of (i-c)
        dp[i] = min(dp[i], dp[i - c] + 1);
      }
    }
  }
  return (dp[target] == INF) ? -1 : dp[target];
}
\end{minted}

\subsection{CoinChange(NumberOfWays)}
\begin{minted}[
    breaklines = true,
    breakanywhere = true,
    frame=lines,    
    fontfamily=tt,
    linenos=false,
    numberblanklines=true,
    numbersep=2pt,
    gobble=0,
    framerule=1pt,
    framesep=1mm,
    funcnamehighlighting=true,
    tabsize=2,
    obeytabs=false,
    mathescape=false
    samepage=false, %with this setting you can force the list to appear on the same page
    showspaces=false,
    showtabs =false,
    texcl=false
]{C++}

ll count_ways(int target, vector<int>& coins) {
  // DP Definition: dp[i] = number of ways to make value i using given coins
  vector<ll> dp(target + 1, 0);
  // Base Case: There is 1 way to make value 0 (by choosing nothing)
  dp[0] = 1;
  for (int c : coins) {  // Loop coins *outside* to count combinations (unordered)
    for (int i = c; i <= target; i++) {
      // Transition: Add ways to make the remainder (i - c)
      dp[i] += dp[i - c];
    }
  }
  return dp[target];
}
\end{minted}

\subsection{CountSubsetWithSumK}
\begin{minted}[
    breaklines = true,
    breakanywhere = true,
    frame=lines,    
    fontfamily=tt,
    linenos=false,
    numberblanklines=true,
    numbersep=2pt,
    gobble=0,
    framerule=1pt,
    framesep=1mm,
    funcnamehighlighting=true,
    tabsize=2,
    obeytabs=false,
    mathescape=false
    samepage=false, %with this setting you can force the list to appear on the same page
    showspaces=false,
    showtabs =false,
    texcl=false
]{C++}

int memo[105][10005];  // (N, Sum)
vector<int> nums;
int K;
int solve(int idx, int sum) {
  // DP: dp[idx][sum] = number of ways to complete the sum using suffix nums[idx...]
  // Base: End of array. Return 1 if sum matches K, else 0
  if (idx == nums.size()) {
    return sum == K;
  }
  if (memo[idx][sum] != -1) return memo[idx][sum];
  // Trans: Sum of ways (Include current element) + (Exclude current element)
  int take = (sum + nums[idx] <= K) ? solve(idx + 1, sum + nums[idx]) : 0;
  int skip = solve(idx + 1, sum);
  return memo[idx][sum] = take + skip;
}
\end{minted}

\subsection{Digit DP}
\begin{minted}[
    breaklines = true,
    breakanywhere = true,
    frame=lines,    
    fontfamily=tt,
    linenos=false,
    numberblanklines=true,
    numbersep=2pt,
    gobble=0,
    framerule=1pt,
    framesep=1mm,
    funcnamehighlighting=true,
    tabsize=2,
    obeytabs=false,
    mathescape=false
    samepage=false, %with this setting you can force the list to appear on the same page
    showspaces=false,
    showtabs =false,
    texcl=false
]{C++}

string S;
ll memo[20][2][180];

ll dp(int idx, bool tight, int current_sum) {
  // DP Definition:
  // dp(idx, tight, current_sum) returns the count of valid suffixes we can form
  // starting from index 'idx', given the 'tight' constraint and the current state (sum).
  // Base Case: We have placed digits for all positions
  if (idx == S.size()) {
    return 1;  // Found 1 valid number (or check: return current_sum == target ? 1 : 0)
  }
  ll &ans = memo[idx][tight][current_sum];
  if (ans != -1) {
    return ans;
  }
  ll ans = 0;
  // Determine the upper bound for the current digit
  int limit = tight ? (S[idx] - '0') : 9;
  for (int digit = 0; digit <= limit; digit++) {
    // Transition:
    // Move to next index (idx + 1).
    // New tight constraint: (tight && (digit == limit))
    // New state: Update current_sum (or other state variable)
    bool next_tight = tight && (digit == limit);
    // Example logic: filtering or constraints go here
    // if (digit == 4) continue; // Example: skip digit 4
    ans += dp(idx + 1, next_tight, current_sum + digit);
  }
  return memo[idx][tight][current_sum] = ans;
}
// Wrapper function to query for range [0, N]
ll query(ll N) {
  if (N < 0) return 0;
  if (N == 0) return 1;  // Handling 0 case if needed
  S = to_string(N);
  memset(memo, -1, sizeof(memo));
  // Start from index 0, tight constraint is initially TRUE, initial sum is 0
  return dp(0, true, 0);
}
// To get answer for range [L, R]: query(R) - query(L - 1)
\end{minted}

\subsection{EditDistance}
\begin{minted}[
    breaklines = true,
    breakanywhere = true,
    frame=lines,    
    fontfamily=tt,
    linenos=false,
    numberblanklines=true,
    numbersep=2pt,
    gobble=0,
    framerule=1pt,
    framesep=1mm,
    funcnamehighlighting=true,
    tabsize=2,
    obeytabs=false,
    mathescape=false
    samepage=false, %with this setting you can force the list to appear on the same page
    showspaces=false,
    showtabs =false,
    texcl=false
]{C++}

int edit_distance(string s1, string s2) {
  int m = s1.size(), n = s2.size();
  // DP: dp[i][j] = min ops to convert s1[0...i-1] to s2[0...j-1]
  vector<vector<int>> dp(m + 1, vector<int>(n + 1));
  // Base: Converting empty string to s2 requires j inserts (and vice versa)
  for (int i = 0; i <= m; i++) dp[i][0] = i;
  for (int j = 0; j <= n; j++) dp[0][j] = j;
  for (int i = 1; i <= m; i++) {
    for (int j = 1; j <= n; j++) {
      if (s1[i - 1] == s2[j - 1]) {
        dp[i][j] = dp[i - 1][j - 1];
      } else {
        // Trans: 1 + min(Delete, Insert, Replace)
        dp[i][j] = 1 + min({dp[i - 1][j], dp[i][j - 1], dp[i - 1][j - 1]});
      }
    }
  }
  return dp[m][n];
}
\end{minted}

\subsection{LIS}
\begin{minted}[
    breaklines = true,
    breakanywhere = true,
    frame=lines,    
    fontfamily=tt,
    linenos=false,
    numberblanklines=true,
    numbersep=2pt,
    gobble=0,
    framerule=1pt,
    framesep=1mm,
    funcnamehighlighting=true,
    tabsize=2,
    obeytabs=false,
    mathescape=false
    samepage=false, %with this setting you can force the list to appear on the same page
    showspaces=false,
    showtabs =false,
    texcl=false
]{C++}

void LIS(const vector<int>& nums) {
  int n = nums.size();
  if (n == 0) return;
  // parent[i] stores the index of the predecessor of nums[i] in the LIS
  vector<int> parent(n, -1);
  // tails_id[k] stores the index (in nums) of the smallest tail of all increasing subsequences of length k+1
  vector<int> tails_id;
  // tails_val[k] stores the actual value of that tail (helper for binary search)
  vector<int> tails_val;
  // DP Definition (Implicit in greedy approach):
  // L[i] = Length of the longest increasing subsequence ending at index i.
  // We don't store L[i] explicitly, but compute it on the fly using the 'tails' arrays.
  // Base Case:
  // The first element starts a subsequence of length 1.
  for (int i = 0; i < n; i++) {
    int x = nums[i];
    // Find the first element in tails_val that is >= x
    auto it = lower_bound(tails_val.begin(), tails_val.end(), x);
    int idx = it - tails_val.begin();
    // Transition:
    // If x extends the longest existing subsequence, append it.
    // Otherwise, replace the existing tail to maintain the "smallest ending value" property
    // which allows future elements to extend the list easier.
    // We also link nums[i] to the element currently ending the subsequence of length (idx).
    if (it == tails_val.end()) {
      tails_val.push_back(x);
      tails_id.push_back(i);
    } else {
      *it = x;
      tails_id[idx] = i;
    }
    // If we are extending a sequence of length > 0, the parent is the end of the previous length
    if (idx > 0) {
      parent[i] = tails_id[idx - 1];
    }
  }
  // Reconstruction
  int curr = tails_id.back();  // Index of the last element of the LIS
  vector<int> lis_path;
  while (curr != -1) {
    lis_path.push_back(nums[curr]);
    curr = parent[curr];
  }
  reverse(lis_path.begin(), lis_path.end());
  // Output
  cout << "LIS Length: " << lis_path.size() << endl;
  for (int x : lis_path) cout << x << " ";
  cout << endl;
}
\end{minted}

\subsection{Maximum Subarray Sum(Kadanes)}
\begin{minted}[
    breaklines = true,
    breakanywhere = true,
    frame=lines,    
    fontfamily=tt,
    linenos=false,
    numberblanklines=true,
    numbersep=2pt,
    gobble=0,
    framerule=1pt,
    framesep=1mm,
    funcnamehighlighting=true,
    tabsize=2,
    obeytabs=false,
    mathescape=false
    samepage=false, %with this setting you can force the list to appear on the same page
    showspaces=false,
    showtabs =false,
    texcl=false
]{C++}

ll max_subarray_sum(const vector<int>& nums) {
  if (nums.empty()) return 0;
  ll current_max = nums[0];
  ll global_max = nums[0];
  for (size_t i = 1; i < nums.size(); i++) {
    current_max = max((ll)nums[i], current_max + nums[i]);
    global_max = max(global_max, current_max);
  }
  return global_max;
}
\end{minted}

\section{Graph}
\subsection{Articulation Point}
\begin{minted}[
    breaklines = true,
    breakanywhere = true,
    frame=lines,    
    fontfamily=tt,
    linenos=false,
    numberblanklines=true,
    numbersep=2pt,
    gobble=0,
    framerule=1pt,
    framesep=1mm,
    funcnamehighlighting=true,
    tabsize=2,
    obeytabs=false,
    mathescape=false
    samepage=false, %with this setting you can force the list to appear on the same page
    showspaces=false,
    showtabs =false,
    texcl=false
]{C++}

int n;                    // number of nodes
vector<vector<int>> lst;  // adjacency list of graph
vector<bool> vis;
vector<int> tin, low;
int timer;
void dfs(int u, int p = -1) {
  vis[u] = true;
  tin[u] = low[u] = timer++;
  int children = 0;
  for (int v : lst[u]) {
    if (v == p) continue;
    if (vis[v]) {
      low[u] = min(low[u], tin[v]);
    } else {
      dfs(v, u);
      low[u] = min(low[u], low[v]);
      if (low[v] >= tin[u] && p != -1){
        IS_CUTPOINT(u);
      }
      ++children;
    }
  }
  // if no vertex below v can reach u or higher → removing u disconnects that subtree
  if (p == -1 && children > 1){
    IS_CUTPOINT(u);
  }
}
void find_cutpoints() {
  timer = 0;
  vis.assign(n, false);
  tin.assign(n, -1);
  low.assign(n, -1);
  for (int i = 0; i < n; ++i) {
    if (!vis[i]){
      dfs(i);
    }
  }
}
\end{minted}

\subsection{BFS}
\begin{minted}[
    breaklines = true,
    breakanywhere = true,
    frame=lines,    
    fontfamily=tt,
    linenos=false,
    numberblanklines=true,
    numbersep=2pt,
    gobble=0,
    framerule=1pt,
    framesep=1mm,
    funcnamehighlighting=true,
    tabsize=2,
    obeytabs=false,
    mathescape=false
    samepage=false, %with this setting you can force the list to appear on the same page
    showspaces=false,
    showtabs =false,
    texcl=false
]{C++}

vector<vector<int>> adj;  // adjacency list representation
int n;                    // number of nodes
int s;                    // source vertex

void bfs() {
  queue<int> q;
  vector<int> d(n), p(n);
  vector<bool> used(n);

  q.push(s);
  used[s] = true;
  p[s] = -1;
  while (!q.empty()) {
    int v = q.front();
    q.pop();
    for (int u : adj[v]) {
      if (!used[u]) {
        used[u] = true;
        q.push(u);
        d[u] = d[v] + 1;
        p[u] = v;
      }
    }
  }
}

// retrieving shortest path

if (!used[u]) {
  cout << "No path!";
} else {
  vector<int> path;
  for (int v = u; v != -1; v = p[v])
    path.push_back(v);
  reverse(path.begin(), path.end());
  cout << "Path: ";
  for (int v : path)
    cout << v << " ";
}
\end{minted}

\subsection{Bellman Ford}
\begin{minted}[
    breaklines = true,
    breakanywhere = true,
    frame=lines,    
    fontfamily=tt,
    linenos=false,
    numberblanklines=true,
    numbersep=2pt,
    gobble=0,
    framerule=1pt,
    framesep=1mm,
    funcnamehighlighting=true,
    tabsize=2,
    obeytabs=false,
    mathescape=false
    samepage=false, %with this setting you can force the list to appear on the same page
    showspaces=false,
    showtabs =false,
    texcl=false
]{C++}

#define ll long long
#define INF 1e18

void solve() {
  int n, m, v;
  cin >> n >> m >> v;  // n = nodes, m = edges, v = source (0-indexed)

  vector<array<ll, 3>> edges(m);  // each edge: {a, b, cost}
  for (int i = 0; i < m; i++) cin >> edges[i][0] >> edges[i][1] >> edges[i][2];

  vector<ll> d(n, INF);
  vector<int> p(n, -1);
  d[v] = 0;

  int x = -1;
  for (int i = 0; i < n; i++) {
    x = -1;
    for (auto& e : edges) {
      int a = e[0], b = e[1];
      ll cost = e[2];
      if (d[a] < INF && d[b] > d[a] + cost) {
        d[b] = max(-INF, d[a] + cost);
        p[b] = a;
        x = b;
      }
    }
  }

  if (x == -1) {
    cout << "No negative cycle from vertex " << v << '\n';
    return;
  }

  int y = x;
  for (int i = 0; i < n; i++) y = p[y];

  vector<int> path;
  for (int cur = y;; cur = p[cur]) {
    path.push_back(cur);
    if (cur == y && path.size() > 1) break;
  }
  reverse(path.begin(), path.end());

  cout << "Negative cycle: ";
  for (int u : path) cout << u << ' ';
  cout << '\n';
}
\end{minted}

\subsection{Bridge Finding DFS}
\begin{minted}[
    breaklines = true,
    breakanywhere = true,
    frame=lines,    
    fontfamily=tt,
    linenos=false,
    numberblanklines=true,
    numbersep=2pt,
    gobble=0,
    framerule=1pt,
    framesep=1mm,
    funcnamehighlighting=true,
    tabsize=2,
    obeytabs=false,
    mathescape=false
    samepage=false, %with this setting you can force the list to appear on the same page
    showspaces=false,
    showtabs =false,
    texcl=false
]{C++}

const int MX = 1e5 + 10;
int n, m, timer = 0;
vector<int> adj[MX];
vector<int> tin(MX, -1), low(MX, -1);
vector<bool> vis(MX, false);
void is_bridge(int u, int v) {
  // do something with the edge
}
void dfs(int u, int p = -1) {
  vis[u] = true;
  tin[u] = low[u] = timer++;
  for (int v : adj[u]) {
    if (v == p) continue;
    if (vis[v]) {
      low[u] = min(low[u], tin[v]);
    } else {
      dfs(v, u);
      low[u] = min(low[u], low[v]);
      if (low[v] > tin[u]) {
        is_bridge(u, v);
      }
    }
  }
}
\end{minted}

\subsection{Cycle Detection in DAG}
\begin{minted}[
    breaklines = true,
    breakanywhere = true,
    frame=lines,    
    fontfamily=tt,
    linenos=false,
    numberblanklines=true,
    numbersep=2pt,
    gobble=0,
    framerule=1pt,
    framesep=1mm,
    funcnamehighlighting=true,
    tabsize=2,
    obeytabs=false,
    mathescape=false
    samepage=false, %with this setting you can force the list to appear on the same page
    showspaces=false,
    showtabs =false,
    texcl=false
]{C++}

const int MX = 1e5 + 10;
bool vis[MX], pathVis[MX];
vector<int> lst[MX];
bool dfs(int u) {
  vis[u] = true;
  pathVis[u] = true;
  for (auto v : lst[u]) {
    if (!vis[v]) {
      if (dfs(v))
        return true;
    } else if (pathVis[v]) {
      return true;
    }
  }
  pathVis[u] = false;
  return false;
}
void solve() {
  // take graph input
  for (int i = 0; i < n; ++i) {
    if (!vis[i])
      dfs(i);
  }
}
\end{minted}

\subsection{DSU}
\begin{minted}[
    breaklines = true,
    breakanywhere = true,
    frame=lines,    
    fontfamily=tt,
    linenos=false,
    numberblanklines=true,
    numbersep=2pt,
    gobble=0,
    framerule=1pt,
    framesep=1mm,
    funcnamehighlighting=true,
    tabsize=2,
    obeytabs=false,
    mathescape=false
    samepage=false, %with this setting you can force the list to appear on the same page
    showspaces=false,
    showtabs =false,
    texcl=false
]{C++}

const int MX = 1e5 + 10;
int par[MX], sz[MX];
void init() {
  for (int i = 1; i < MX; i++) {
    par[i] = i;
    sz[i] = 1;
  }
}
int findpar(int x) {
  if (par[x] == x) return x;
  return par[x] = findpar(par[x]);
}
void unite(int u, int v) {
  u = findpar(u);
  v = findpar(v);
  if (u != v) {
    if (sz[u] < sz[v]) {
      swap(u, v);
    }
    sz[u] += sz[v];
    par[v] = u;
  }
}
\end{minted}

\subsection{Dijkastra}
\begin{minted}[
    breaklines = true,
    breakanywhere = true,
    frame=lines,    
    fontfamily=tt,
    linenos=false,
    numberblanklines=true,
    numbersep=2pt,
    gobble=0,
    framerule=1pt,
    framesep=1mm,
    funcnamehighlighting=true,
    tabsize=2,
    obeytabs=false,
    mathescape=false
    samepage=false, %with this setting you can force the list to appear on the same page
    showspaces=false,
    showtabs =false,
    texcl=false
]{C++}

const int N = 1e5 + 5, INF = 1e18 + 7;
vector<pair<int, int>> g[N];
bool visited[N];
vector<int> dist(N, INF), parent(N);
bool dijkstra(int source) {
  priority_queue<pair<int, int>, vector<pair<int, int>>, greater<pair<int, int>>> pq;
  pq.push({0, source});
  dist[source] = 0;
  parent[source] = -1;
  while (pq.size()) {
    int x = pq.top().second;
    pq.pop();
    if (visited[x]) continue;
    visited[x] = 1;
    for (auto [child_x, child_wt] : g[x]) {
      if (dist[x] + child_wt < dist[child_x]) {
        parent[child_x] = x;
        dist[child_x] = child_wt + dist[x];
        pq.push({dist[child_x], child_x});
      }
    }
  }
  return (dist[n] == INF);
}
\end{minted}

\subsection{Euler Tour}
\begin{minted}[
    breaklines = true,
    breakanywhere = true,
    frame=lines,    
    fontfamily=tt,
    linenos=false,
    numberblanklines=true,
    numbersep=2pt,
    gobble=0,
    framerule=1pt,
    framesep=1mm,
    funcnamehighlighting=true,
    tabsize=2,
    obeytabs=false,
    mathescape=false
    samepage=false, %with this setting you can force the list to appear on the same page
    showspaces=false,
    showtabs =false,
    texcl=false
]{C++}

const int MX = 2e5 + 10;
int timer = -1;
// s = start pos, e = end pos
int val[MX], s[MX], e[MX], flat[MX];
vector<int> lst[MX];
void dfs(int u, int p) {
  s[u] = ++timer;
  flat[timer] = val[u];
  for (auto v : lst[u]) {
    if (v != p)
      dfs(v, u);
  }
  e[u] = timer;
}
\end{minted}

\subsection{Floyd Warshall}
\begin{minted}[
    breaklines = true,
    breakanywhere = true,
    frame=lines,    
    fontfamily=tt,
    linenos=false,
    numberblanklines=true,
    numbersep=2pt,
    gobble=0,
    framerule=1pt,
    framesep=1mm,
    funcnamehighlighting=true,
    tabsize=2,
    obeytabs=false,
    mathescape=false
    samepage=false, %with this setting you can force the list to appear on the same page
    showspaces=false,
    showtabs =false,
    texcl=false
]{C++}

vector<vector<int>> d(n, vector<int>(n, INF));
// take graph input into d
for (int k = 0; k < n; ++k) {
  for (int i = 0; i < n; ++i) {
    for (int j = 0; j < n; ++j) {
      if (d[i][k] < INF && d[k][j] < INF)
        d[i][j] = min(d[i][j], d[i][k] + d[k][j]);
    }
  }
}
\end{minted}

\subsection{MST}
\begin{minted}[
    breaklines = true,
    breakanywhere = true,
    frame=lines,    
    fontfamily=tt,
    linenos=false,
    numberblanklines=true,
    numbersep=2pt,
    gobble=0,
    framerule=1pt,
    framesep=1mm,
    funcnamehighlighting=true,
    tabsize=2,
    obeytabs=false,
    mathescape=false
    samepage=false, %with this setting you can force the list to appear on the same page
    showspaces=false,
    showtabs =false,
    texcl=false
]{C++}

// DSU first
void solve() {
  int n, m;
  cin >> n >> m;
  vector<tuple<int, int, int>> edges;
  for (int i = 0; i < m; ++i) {
    int u, v, wt;
    cin >> u >> v >> wt;
    edges.push_back({wt, u, v});
  }
  sort(edges.begin(), edges.end());
  init(n);
  int cost = 0;
  for (tuple& [ wt, u, v ] : edges) {
    if (findpar(u) == findpar(v)) continue;
    unite(u, v);
    cost += wt;
  }
  cout << cost << endl;
}
\end{minted}

\subsection{Max Bipartite Matching[Hopcroft Karp]}
\begin{minted}[
    breaklines = true,
    breakanywhere = true,
    frame=lines,    
    fontfamily=tt,
    linenos=false,
    numberblanklines=true,
    numbersep=2pt,
    gobble=0,
    framerule=1pt,
    framesep=1mm,
    funcnamehighlighting=true,
    tabsize=2,
    obeytabs=false,
    mathescape=false
    samepage=false, %with this setting you can force the list to appear on the same page
    showspaces=false,
    showtabs =false,
    texcl=false
]{C++}

const int INF = 1e9;
void hopcroftCarp() {
  int n, m, e;
  cin >> n >> m >> e;
  vector<int> adj[n];
  for (int i = 0; i < e; ++i) {
    int u, v;
    cin >> u >> v;
    --u;
    --v;
    adj[u].push_back(v);
  }
  vector<int> ml(m, -1), mr(n, -1), dist(n);
  auto bfs = [&]() -> bool {
    queue<int> q;
    for (int u = 0; u < n; ++u) {
      if (mr[u] == -1) {
        dist[u] = 0;
        q.push(u);
      } else {
        dist[u] = INF;
      }
    }
    bool foundAugmenting = false;
    while (!q.empty()) {
      int u = q.front();
      q.pop();
      for (int v : adj[u]) {
        int pairedLeft = ml[v];
        if (pairedLeft == -1) {
          foundAugmenting = true;
        } else if (dist[pairedLeft] == INF) {
          dist[pairedLeft] = dist[u] + 1;
          q.push(pairedLeft);
        }
      }
    }
    return foundAugmenting;
  };
  function<bool(int)> dfs = [&](int u) -> bool {
    for (int v : adj[u]) {
      int pairedLeft = ml[v];
      if (pairedLeft == -1 or (dist[pairedLeft] == dist[u] + 1 and dfs(pairedLeft))) {
        mr[u] = v;
        ml[v] = u;
        return true;
      }
    }
    dist[u] = INF;
    return false;
  };
  int matching = 0;
  while (bfs()) {
    for (int u = 0; u < n; ++u) {
      if (mr[u] == -1) {
        if (dfs(u)) matching++;
      }
    }
  }
  cout << matching << el;
  for (int u = 0; u < n; ++u) {
    if (mr[u] != -1) {
      cout << u << " " << mr[u] << el;
    }
  }
}
\end{minted}

\subsection{Max Bipartite Matching[Kuhn's]}
\begin{minted}[
    breaklines = true,
    breakanywhere = true,
    frame=lines,    
    fontfamily=tt,
    linenos=false,
    numberblanklines=true,
    numbersep=2pt,
    gobble=0,
    framerule=1pt,
    framesep=1mm,
    funcnamehighlighting=true,
    tabsize=2,
    obeytabs=false,
    mathescape=false
    samepage=false, %with this setting you can force the list to appear on the same page
    showspaces=false,
    showtabs =false,
    texcl=false
]{C++}

// left set size, right set size, edge count
int n, k, m, visToken = 1;
vector<int> lst[MX];
int mr[MX], ml[MX], vis[MX];
bool try_kuhn(int u) {
  if (vis[u] == visToken)
    return false;
  vis[u] = visToken;
  for (auto v : lst[u]) {
    if (ml[v] == -1 or try_kuhn(ml[v])) {
      ml[v] = u;
      mr[u] = v;
      return true;
    }
  }
  return false;
}
void solve() {
  cin >> n >> k >> m;
  for (int i = 0; i < m; ++i) {
    int u, v;
    cin >> u >> v;
    --u, --v;
    lst[u].push_back(v);
  }
  fill(mr, mr + n, -1);
  fill(ml, ml + k, -1);
  int ans = 0;
  for (int u = 0; u < n; ++u) {
    for (auto v : lst[u]) {
      if (ml[v] == -1) {
        ml[v] = u;
        mr[u] = v;
        ans++;
        break;
      }
    }
  }
  for (int u = 0; u < n; ++u) {
    if (mr[u] != -1) continue;
    visToken++;
    if (try_kuhn(u))
      ans++;
  }
  cout << ans << el;
  for (int v = 0; v < k; ++v) {
    if (ml[v] != -1) {
      cout << ml[v] + 1 << " " << v + 1 << el;
    }
  }
}
\end{minted}

\subsection{Max Flow}
\begin{minted}[
    breaklines = true,
    breakanywhere = true,
    frame=lines,    
    fontfamily=tt,
    linenos=false,
    numberblanklines=true,
    numbersep=2pt,
    gobble=0,
    framerule=1pt,
    framesep=1mm,
    funcnamehighlighting=true,
    tabsize=2,
    obeytabs=false,
    mathescape=false
    samepage=false, %with this setting you can force the list to appear on the same page
    showspaces=false,
    showtabs =false,
    texcl=false
]{C++}

struct Dinic {
  struct Edge {
    int to;
    ll capacity;
    int rev;  // index of reverse edge
  };

  vector<vector<Edge>> adj;
  vector<int> level;
  vector<int> ptr;
  int n;

  Dinic(int nodes) : n(nodes), adj(nodes), level(nodes), ptr(nodes) {}

  void add_edge(int from, int to, ll cap) {
    adj[from].push_back({to, cap, (int)adj[to].size()});
    adj[to].push_back({from, 0, (int)adj[from].size() - 1});
  }

  bool bfs(int s, int t) {
    fill(level.begin(), level.end(), -1);
    level[s] = 0;
    queue<int> q;
    q.push(s);
    while (!q.empty()) {
      int v = q.front();
      q.pop();
      for (auto& edge : adj[v]) {
        if (edge.capacity > 0 && level[edge.to] == -1) {
          level[edge.to] = level[v] + 1;
          q.push(edge.to);
        }
      }
    }
    return level[t] != -1;
  }

  ll dfs(int v, int t, ll pushed) {
    if (pushed == 0) return 0;
    if (v == t) return pushed;
    for (int& cid = ptr[v]; cid < adj[v].size(); ++cid) {
      auto& edge = adj[v][cid];
      int tr = edge.to;
      if (level[v] + 1 != level[tr] || edge.capacity == 0) continue;
      ll tr_pushed = dfs(tr, t, min(pushed, edge.capacity));
      if (tr_pushed == 0) continue;
      edge.capacity -= tr_pushed;
      adj[tr][edge.rev].capacity += tr_pushed;
      return tr_pushed;
    }
    return 0;
  }

  ll max_flow(int s, int t) {
    ll flow = 0;
    while (bfs(s, t)) {
      fill(ptr.begin(), ptr.end(), 0);
      while (ll pushed = dfs(s, t, 1e18)) {
        flow += pushed;
      }
    }
    return flow;
  }
};
// Usage in int main():
int n, m;
cin >> n >> m;
Dinic dinic(n + 1); // Initialize with N+1 nodes (0 to N)
for (int i = 0; i < m; i++) {
    int u, v; ll cap;
    cin >> u >> v >> cap;
    // Add directed edge u -> v with capacity
    // For bidirectional, add edge v -> u as well
    dinic.add_edge(u, v, cap);
}
// Compute flow from node 1 (Source) to node n (Sink)
cout << dinic.max_flow(1, n) << endl;
\end{minted}

\subsection{Tarjan's Algorithm}
\begin{minted}[
    breaklines = true,
    breakanywhere = true,
    frame=lines,    
    fontfamily=tt,
    linenos=false,
    numberblanklines=true,
    numbersep=2pt,
    gobble=0,
    framerule=1pt,
    framesep=1mm,
    funcnamehighlighting=true,
    tabsize=2,
    obeytabs=false,
    mathescape=false
    samepage=false, %with this setting you can force the list to appear on the same page
    showspaces=false,
    showtabs =false,
    texcl=false
]{C++}

#include <bits/stdc++.h>
#define int long long
#define endl '\n'
using namespace std;

const int N = 55;
vector<int> g[N];
vector<int>timee(N,0),low(N,0);
int vis[N];
int n,m;
int cnt = 0;
int timer = 1;
void dfs(int vertex, int par = -1) {
    vis[vertex] = 1;
    timee[vertex] = low[vertex] = timer;
    ++timer;
    for(auto child : g[vertex]) {
      if(child == par ) continue;
      if(vis[child] == 1 ) {
          low[vertex] = min(low[vertex],low[child]); continue;
      }
      dfs(child, vertex);
      low[vertex] = min(low[vertex],low[child]);
      if(low[child] > timee[vertex]) ++cnt;
    }
}
void solve()
{
    cin >> n >> m;
    for(int i = 0; i < m; i++) {
      int x,y; cin >> x >> y;
      g[x].push_back(y);
      g[y].push_back(x);
    }
    dfs(1);
    cout << cnt << endl;
}
signed main()
{
    ios_base::sync_with_stdio(0),cin.tie(0);
    int tt = 1;
   // cin >> tt;
    while(tt--)
    solve();
}
\end{minted}

\subsection{Topological Sorting}
\begin{minted}[
    breaklines = true,
    breakanywhere = true,
    frame=lines,    
    fontfamily=tt,
    linenos=false,
    numberblanklines=true,
    numbersep=2pt,
    gobble=0,
    framerule=1pt,
    framesep=1mm,
    funcnamehighlighting=true,
    tabsize=2,
    obeytabs=false,
    mathescape=false
    samepage=false, %with this setting you can force the list to appear on the same page
    showspaces=false,
    showtabs =false,
    texcl=false
]{C++}

const int N = 1e5 + 10;
vector<int> g[N], indegree(N, 0);
vector<int> topSort(int n) {
  queue<int> q;
  vector<int> order;
  for (int i = 1; i <= n; i++) {
    if (indegree[i] == 0) {
      q.push(i);
    }
  }
  while (!q.empty()) {
    int u = q.front();
    q.pop();
    order.push_back(u);
    for (int v : g[u]) {
      indegree[v]--;
      if (indegree[v] == 0) {
        q.push(v);
      }
    }
  }
  return order;
}
\end{minted}

\subsection{Weighted Union Find}
\begin{minted}[
    breaklines = true,
    breakanywhere = true,
    frame=lines,    
    fontfamily=tt,
    linenos=false,
    numberblanklines=true,
    numbersep=2pt,
    gobble=0,
    framerule=1pt,
    framesep=1mm,
    funcnamehighlighting=true,
    tabsize=2,
    obeytabs=false,
    mathescape=false
    samepage=false, %with this setting you can force the list to appear on the same page
    showspaces=false,
    showtabs =false,
    texcl=false
]{C++}

const int MX = 2e5 + 10;
int par[MX], sz[MX];
ll d[MX];
void init() {
  for (int i = 0; i < MX; ++i) {
    par[i] = i;
    sz[i] = 1;
    d[i] = 0;
  }
}
int findpar(int x) {
  if (par[x] == x) return x;
  int p = par[x];
  par[x] = findpar(p);
  d[x] += d[p];
  return par[x];
}
bool unite(int a, int b, ll w) {
  int ra = findpar(a);
  int rb = findpar(b);
  if (ra == rb) {
    return (d[b] - d[a] == w);
  }
  if (sz[ra] < sz[rb]) {
    swap(a, b);
    swap(ra, rb);
    w = -w;
  }
  par[rb] = ra;
  d[rb] = d[a] + w - d[b];
  sz[ra] += sz[rb];
  return true;
}
ll dist(int a, int b) {
  findpar(a), findpar(b);
  return d[b] - d[a];
}
\end{minted}

\section{Math}
\subsection{ExtendedEuclid}
\begin{minted}[
    breaklines = true,
    breakanywhere = true,
    frame=lines,    
    fontfamily=tt,
    linenos=false,
    numberblanklines=true,
    numbersep=2pt,
    gobble=0,
    framerule=1pt,
    framesep=1mm,
    funcnamehighlighting=true,
    tabsize=2,
    obeytabs=false,
    mathescape=false
    samepage=false, %with this setting you can force the list to appear on the same page
    showspaces=false,
    showtabs =false,
    texcl=false
]{C++}

// It finds integers x and y (coefficients) such that: ax + by = gcd(a, b)
int extendedGCD(int a, int b, int& x, int& y) {
  // Base Case
  if (a == 0) {
    x = 0;
    y = 1;
    return b;
  }
  int x1, y1;
  int g = extendedGCD(b % a, a, x1, y1);
  // Update x and y using results of recursive call
  x = y1 - (b / a) * x1;
  y = x1;
  return g;
}
\end{minted}

\subsection{Formula}
\begin{minted}[
    breaklines = true,
    breakanywhere = true,
    frame=lines,    
    fontfamily=tt,
    linenos=false,
    numberblanklines=true,
    numbersep=2pt,
    gobble=0,
    framerule=1pt,
    framesep=1mm,
    funcnamehighlighting=true,
    tabsize=2,
    obeytabs=false,
    mathescape=false
    samepage=false, %with this setting you can force the list to appear on the same page
    showspaces=false,
    showtabs =false,
    texcl=false
]{C++}

/*
1. Sum of Pair Sums
(n - 1) * Σ a[k]

2. Sum of Subarray Sums
Σ ( a[k] * k * (n - k + 1) )

3. Sum of Subset Sums
2^(n - 1) * Σ a[k]

4. Product of Pair Products
Π ( a[k]^(n - 1) )

5. XOR of Subarray XORs
⊕ { a[k] | k(n - k + 1) is odd }

6. Sum of Max - Min over all Subsets
Σ ( a[k] * (2^(k - 1) - 2^(n - k)) )

7. Sum of Sum*Length over all Subarrays
Σ ( a[k] * ( k(n - k + 1)(n + 1) / 2 ) )

8. Sum of Pair XORs
Σ ( cnt1(b) * cnt0(b) * 2^b )

9. Sum of Pair ANDs
Σ ( C(cnt1(b), 2) * 2^b )

10. Sum of Pair ORs
Σ ( (C(n, 2) - C(cnt0(b), 2)) * 2^b )

11. Sum of Subset XORs
Σ ( [cnt1(b) > 0] * 2^(n - 1) * 2^b )

12. Sum of Subset ANDs
Σ ( (2^(cnt1(b)) - 1) * 2^b )

13. Sum of Subset ORs
Σ ( (2^n - 2^(cnt0(b))) * 2^b )
*/
\end{minted}

\subsection{Matrix Exponentiation}
\begin{minted}[
    breaklines = true,
    breakanywhere = true,
    frame=lines,    
    fontfamily=tt,
    linenos=false,
    numberblanklines=true,
    numbersep=2pt,
    gobble=0,
    framerule=1pt,
    framesep=1mm,
    funcnamehighlighting=true,
    tabsize=2,
    obeytabs=false,
    mathescape=false
    samepage=false, %with this setting you can force the list to appear on the same page
    showspaces=false,
    showtabs =false,
    texcl=false
]{C++}

const ll mod = 1e9;
vector<vector<ll>> matMul(vector<vector<ll>>& a, vector<vector<ll>>& b) {
  ll row1 = a.size(), col1 = a[0].size();
  ll row2 = b.size(), col2 = b[0].size();
  vector<vector<ll>> res(row1, vector<ll>(col2, 0));
  for (ll i = 0; i < row1; i++) {
    for (ll j = 0; j < col1; j++) {
      for (ll k = 0; k < row2; k++) {
        res[i][j] = (res[i][j] + (1LL * a[i][k] * b[k][j]) % mod) % mod;
      }
    }
  }
  return res;
}
vector<vector<ll>> matExpo(vector<vector<ll>>& Mat, ll exp) {
  ll row = Mat.size(), col = Mat[0].size();
  ll p = row;
  vector<vector<ll>> res(p, vector<ll>(p, 0));
  for (ll i = 0; i < p; i++) res[i][i] = 1;
  while (exp) {
    if (exp & 1) res = matMul(res, Mat);
    Mat = matMul(Mat, Mat);
    exp >>= 1;
  }
  return res;
}
// b = (A(i), A(i-1), A(i-2), A(i-3))
// M = Magic matrix, nth = nth term, known = known value
ll get_nth(ll nth, ll known, vector<ll>& b, vector<vector<ll>>& M) {
  if (nth <= known) return b[nth - 1] % mod;
  reverse(b.begin(), b.end());
  vector<vector<ll>> me = matExpo(M, nth - known);  // MAT^(nth-known)
  ll ans = 0;
  for (int i = 0; i < known; i++) {
    ans = (ans + (b[i] * me[i][0]) % mod) % mod;
  }
  return ans;
}
\end{minted}

\subsection{Matrix Rotation}
\begin{minted}[
    breaklines = true,
    breakanywhere = true,
    frame=lines,    
    fontfamily=tt,
    linenos=false,
    numberblanklines=true,
    numbersep=2pt,
    gobble=0,
    framerule=1pt,
    framesep=1mm,
    funcnamehighlighting=true,
    tabsize=2,
    obeytabs=false,
    mathescape=false
    samepage=false, %with this setting you can force the list to appear on the same page
    showspaces=false,
    showtabs =false,
    texcl=false
]{C++}

//90* clock-wise
now = {{0, 1, 0}, { -1, 0, 0}, {0, 0, 1}};
//90* anti-clock
now = {{0, -1, 0}, {1, 0, 0}, {0, 0, 1}};
//mirror with x axis at point p
now = {{ -1, 0, 2 * p}, {0, 1, 0}, {0, 0, 1}};
//mirror with y axis at point p
now = {{1, 0, 0}, {0, -1, 2 * p}, {0, 0, 1}};
op[i + 1] = matMul(now, op[i]); // this
// op[i + 1] = matMul(op[i], now); //not this
\end{minted}

\section{Miscellaneous}
\subsection{Max Subarray Size Sum equal K}
\begin{minted}[
    breaklines = true,
    breakanywhere = true,
    frame=lines,    
    fontfamily=tt,
    linenos=false,
    numberblanklines=true,
    numbersep=2pt,
    gobble=0,
    framerule=1pt,
    framesep=1mm,
    funcnamehighlighting=true,
    tabsize=2,
    obeytabs=false,
    mathescape=false
    samepage=false, %with this setting you can force the list to appear on the same page
    showspaces=false,
    showtabs =false,
    texcl=false
]{C++}

//write gpHashTable code before this part
void solution(){
int n, k; cin >> n >> k;
int total_sum = 0;
vector < int > pre(n + 7, 0);
for (int i = 1; i <= n; i++) {
  int temp; cin >> temp;
  total_sum += temp;
  if (i == 1) pre[i] = temp;
  else pre[i] = pre[i - 1] + temp;
}
if (total_sum < k) {
  cout << "-1" << endl; return;
}
if (total_sum == k) {
  cout << "0" << endl; return;
}
int maximum_subSize = 0;
gp_hash_table < int, int, customHash> table;
for (int i = 1; i <= n; i++) {
  if (pre[i] >= k) {
    int subSUM = pre[i] - k;
    if (subSUM == 0){
      maximum_subSize = max(maximum_subSize, i);
    }
    else if (table[subSUM]) {
      int left = table[subSUM];
      int right = i; int subSize = right - left;
      maximum_subSize = max(subSize, maximum_subSize);
    }
  } if (!table[pre[i]]) table[pre[i]] = i;
} cout << maximum_subSize << endl; }
\end{minted}

\subsection{Merge Sort}
\begin{minted}[
    breaklines = true,
    breakanywhere = true,
    frame=lines,    
    fontfamily=tt,
    linenos=false,
    numberblanklines=true,
    numbersep=2pt,
    gobble=0,
    framerule=1pt,
    framesep=1mm,
    funcnamehighlighting=true,
    tabsize=2,
    obeytabs=false,
    mathescape=false
    samepage=false, %with this setting you can force the list to appear on the same page
    showspaces=false,
    showtabs =false,
    texcl=false
]{C++}

// use array of elements, if multiple testcase make inv = 0 each time
// int inv = 0;
void merge(int vct[], int l, int m, int r) {
  int left = m - l + 1, right = r - m, lv[left], rv[right];
  for (int i = 0; i < left; i++) {
    lv[i] = vct[l + i];
  }
  for (int i = 0; i < right; i++) {
    rv[i] = vct[m + 1 + i];
  }
  int i = 0, j = 0, to = l;
  while (i < left && j < right) {
    if (lv[i] <= rv[j]) {
      vct[to] = lv[i];
      i++;
    } else {
      vct[to] = rv[j];
      j++;
      // inversion count
      // int pore = left-i; inv+=pore;
    }
    to++;
  }
  while (i < left) {
    vct[to] = lv[i];
    i++;
    to++;
  }
  while (j < right) {
    vct[to] = rv[j];
    j++;
    to++;
  }
}
void merge_sort(int vct[], int l, int r) {
  if (r <= l) return;
  int m = l + ((r - l) / 2);
  merge_sort(vct, l, m);
  merge_sort(vct, m + 1, r);
  merge(vct, l, m, r);
}
\end{minted}

\subsection{Number of Subarray Sum is K}
\begin{minted}[
    breaklines = true,
    breakanywhere = true,
    frame=lines,    
    fontfamily=tt,
    linenos=false,
    numberblanklines=true,
    numbersep=2pt,
    gobble=0,
    framerule=1pt,
    framesep=1mm,
    funcnamehighlighting=true,
    tabsize=2,
    obeytabs=false,
    mathescape=false
    samepage=false, %with this setting you can force the list to appear on the same page
    showspaces=false,
    showtabs =false,
    texcl=false
]{C++}

//write gpHashTable code before this part
void solution(){
int n, k; cin >> n >> k;
int total_sum = 0;
vector < int > pre(n + 7, 0);
for (int i = 1; i <= n; i++) {
  int temp; cin >> temp;
  total_sum += temp;
  if (i == 1) pre[i] = temp;
  else pre[i] = pre[i - 1] + temp; }
int cnt_subarry = 0;
gp_hash_table < int, int, customHash> table;
table[0] = 1;
for (int i = 1; i <= n; i++) {
  cnt_subarry += table[pre[i] - k];
  table[pre[i]]++;
} cout << cnt_subarry << endl; }
\end{minted}

\section{Number Theory}
\subsection{All In One NT}
\begin{minted}[
    breaklines = true,
    breakanywhere = true,
    frame=lines,    
    fontfamily=tt,
    linenos=false,
    numberblanklines=true,
    numbersep=2pt,
    gobble=0,
    framerule=1pt,
    framesep=1mm,
    funcnamehighlighting=true,
    tabsize=2,
    obeytabs=false,
    mathescape=false
    samepage=false, %with this setting you can force the list to appear on the same page
    showspaces=false,
    showtabs =false,
    texcl=false
]{C++}

const int MAXN = 1e6 + 9;
typedef struct info {
  int lowest_prime = 0, greatest_prime = 0, distinct_prime = 0;
  int total_prime = 0, NOD = 0, SOD = 0;
} info;
info num[MAXN];
void preStore() {
  for (int i = 2; i < MAXN; i++) {
    int n = i;
    map<int, int> factors;  // Key->Factor, Val->count
    int SOD = 1, NOD = 1, total_p_factor = 0;
    if (n % 2 == 0) {
      while (n % 2 == 0) {
        n /= 2;
        factors[2]++;
        total_p_factor++;
      }
      SOD *= (1 << (factors[2] + 1)) - 1;
      NOD *= (factors[2] + 1);
    }
    for (int i = 3; i * i <= n; i += 2) {
      if (n % i == 0) {
        while (n % i == 0) {
          n /= i;
          factors[i]++;
          total_p_factor++;
        }
        SOD *= (pow(i, factors[i] + 1) - 1) / (i - 1);
        NOD *= (factors[i] + 1);
      }
    }
    if (n > 1) {
      factors[n]++;
      SOD *= (pow(n, 2) - 1) / (n - 1);
      NOD *= 2;
      total_p_factor++;
    }
    num[i].distinct_prime = factors.size();
    num[i].total_prime = total_p_factor;
    num[i].NOD = NOD;
    num[i].SOD = SOD;
    auto lowest_prime = factors.begin();
    auto greatest_prime = factors.rbegin();
    num[i].lowest_prime = lowest_prime->first;
    num[i].greatest_prime = greatest_prime->first;
  }
}
\end{minted}

\subsection{Divisor Sieve}
\begin{minted}[
    breaklines = true,
    breakanywhere = true,
    frame=lines,    
    fontfamily=tt,
    linenos=false,
    numberblanklines=true,
    numbersep=2pt,
    gobble=0,
    framerule=1pt,
    framesep=1mm,
    funcnamehighlighting=true,
    tabsize=2,
    obeytabs=false,
    mathescape=false
    samepage=false, %with this setting you can force the list to appear on the same page
    showspaces=false,
    showtabs =false,
    texcl=false
]{C++}

const int mxN = 1e5 + 10;
vector<int> divisors[mxN];
void divisorSeive() {
  for (int i = 1; i < mxN; i++) {
    for (int j = i; j < mxN; j += i) {
      divisors[j].push_back(i);
    }
  }
}
\end{minted}

\subsection{Mobius Function}
\begin{minted}[
    breaklines = true,
    breakanywhere = true,
    frame=lines,    
    fontfamily=tt,
    linenos=false,
    numberblanklines=true,
    numbersep=2pt,
    gobble=0,
    framerule=1pt,
    framesep=1mm,
    funcnamehighlighting=true,
    tabsize=2,
    obeytabs=false,
    mathescape=false
    samepage=false, %with this setting you can force the list to appear on the same page
    showspaces=false,
    showtabs =false,
    texcl=false
]{C++}

//FULL MÖBIUS FUNCTION TEMPLATE
//Möbius Sieve (O(N log log N))
//Computes:
//mu[n] → Möbius value
//primes[] → list of primes
//isPrime[]
#include <bits/stdc++.h>
using namespace std;

const int MAXN = 1000000;
int mu[MAXN + 1];
bool isPrime[MAXN + 1];
vector<int> primes;

void mobius_sieve() {
    for (int i = 1; i <= MAXN; i++)
      mu[i] = 1, isPrime[i] = true;
    for (int i = 2; i <= MAXN; i++) {
        if (isPrime[i]) {
            primes.push_back(i);
            for (int j = i; j <= MAXN; j += i) {
                isPrime[j] = false;
                mu[j] *= -1;
            }
            long long sq = 1LL * i * i;
            if (sq <= MAXN) {
                for (int j = sq; j <= MAXN; j += sq)
                  mu[j] = 0;
            }
      }
    }
    mu[1] = 1;
}

//Count Numbers ≤ n Coprime with x
long long count_coprime(long long n, int x) {
    long long ans = 0;
    for (int d = 1; d * d <= x; d++) {
        if (x % d == 0) {
            ans += 1LL * mu[d] * (n / d);
            if (d != x / d)
               ans += 1LL * mu[x / d] * (n / (x / d));
      }
    }
    return ans;
}
//Count Pairs (i, j) with gcd(i, j) = 1 (1 ≤ i,j ≤ n)

long long count_coprime_pairs(int n) {
    long long ans = 0;
    for (int d = 1; d <= n; d++) {
      long long t = n / d;
      ans += 1LL * mu[d] * t * t;
    }
    return ans;
}


//Count Coprime Pairs in an Array (ICPC CLASSIC)
//Problem:
//Count (i,j) such that gcd(a[i], a[j]) = 1.
long long count_array_coprime_pairs(vector<int>& a) {
    int maxA = *max_element(a.begin(), a.end());
    vector<int> cnt(maxA + 1, 0);

    for (int x : a)
      cnt[x]++;

    vector<int> freq(maxA + 1, 0);
    for (int i = 1; i <= maxA; i++) {
      for (int j = i; j <= maxA; j += i)
      freq[i] += cnt[j];
    }
    long long ans = 0;
    for (int d = 1; d <= maxA; d++) {
      if (mu[d] != 0)
          ans += 1LL * mu[d] * freq[d] * (freq[d] - 1) / 2;
    }
    return ans;
}

//Mobius Inversion (Template)

vector<long long> mobius_inversion(vector<long long>& f, int n) {
    vector<long long> g(n + 1, 0);
    for (int i = 1; i <= n; i++) {
      for (int j = i; j <= n; j += i) {
        g[j] += mu[i] * f[j / i];
      }
    }
    return g;
}

//When to Use This Template
//Use Möbius when problem mentions:
//gcd = 1
//coprime pairs
//divisor sums
//inclusion–exclusion
//multiplicative functions

//Complexity Summary
//Task
//Time
//Sieve
//O(N log log N)
//Coprime count
//O(√x)
//Pair counting
//O(N)
//Array coprime
//O(A log A)


//Contest Tips (IMPORTANT)
//Always precompute μ once
//Combine with prefix sums
//Watch out for overflow
//μ values are only {-1, 0, +1}
\end{minted}

\subsection{Number of Pairs with GCD equal g}
\begin{minted}[
    breaklines = true,
    breakanywhere = true,
    frame=lines,    
    fontfamily=tt,
    linenos=false,
    numberblanklines=true,
    numbersep=2pt,
    gobble=0,
    framerule=1pt,
    framesep=1mm,
    funcnamehighlighting=true,
    tabsize=2,
    obeytabs=false,
    mathescape=false
    samepage=false, %with this setting you can force the list to appear on the same page
    showspaces=false,
    showtabs =false,
    texcl=false
]{C++}

/*a[i] <= 1e6
for all 1<=g<=n, how many pairs exist such that g = gcd(a[i], a[j]);
complexity : nlogn
*/
ll n; cin >> n;
ll a[n + 1];
ll cnt[n + 1]; memset(cnt, 0, sizeof cnt);
for (ll i = 1; i <= n; i++) {cin >> a[i]; cnt[a[i]]++;}
ll gcd[n + 1]; memset(gcd, 0, sizeof gcd);
for (ll i = n; i >= 1; i--) {
	ll pair = 0, invalid_pair = 0;
	for (ll j = i; j <= n; j += i) {
		pair += cnt[j];
		invalid_pair += gcd[j];}
	pair = (pair * (pair - 1)) / 2;
	gcd[i] = pair - invalid_pair;
	// how many pairs exist whose gcd is i
}
\end{minted}

\subsection{Phi(1toN)}
\begin{minted}[
    breaklines = true,
    breakanywhere = true,
    frame=lines,    
    fontfamily=tt,
    linenos=false,
    numberblanklines=true,
    numbersep=2pt,
    gobble=0,
    framerule=1pt,
    framesep=1mm,
    funcnamehighlighting=true,
    tabsize=2,
    obeytabs=false,
    mathescape=false
    samepage=false, %with this setting you can force the list to appear on the same page
    showspaces=false,
    showtabs =false,
    texcl=false
]{C++}

const int mxN = 1e7+10;
vector<int> phi(mxN);
void phi_till() { //O(n.log.log(n))
  for (int i = 0; i < mxN; i++) phi[i] = i;
  for (int i = 2; i < mxN; i++) {
    if (phi[i] == i) {
      for (int j = i; j < mxN; j += i){
        phi[j] -= phi[j] / i;
      }
    }
  }
}
\end{minted}

\subsection{Phi}
\begin{minted}[
    breaklines = true,
    breakanywhere = true,
    frame=lines,    
    fontfamily=tt,
    linenos=false,
    numberblanklines=true,
    numbersep=2pt,
    gobble=0,
    framerule=1pt,
    framesep=1mm,
    funcnamehighlighting=true,
    tabsize=2,
    obeytabs=false,
    mathescape=false
    samepage=false, %with this setting you can force the list to appear on the same page
    showspaces=false,
    showtabs =false,
    texcl=false
]{C++}

int phi(int n) {  // sqrt(n)
  int result = n;
  for (int i = 2; i * i <= n; i++) {
    if (n % i == 0) {
      while (n % i == 0) n /= i;
      result -= result / i;
    }
  }
  if (n > 1) result -= result / n;
  return result;
}
\end{minted}

\subsection{SOD NOD}
\begin{minted}[
    breaklines = true,
    breakanywhere = true,
    frame=lines,    
    fontfamily=tt,
    linenos=false,
    numberblanklines=true,
    numbersep=2pt,
    gobble=0,
    framerule=1pt,
    framesep=1mm,
    funcnamehighlighting=true,
    tabsize=2,
    obeytabs=false,
    mathescape=false
    samepage=false, %with this setting you can force the list to appear on the same page
    showspaces=false,
    showtabs =false,
    texcl=false
]{C++}

// SOD = ((P^(x+1)-1)/(P-1)) * ((Q^(y+1)-1)/(Q-1)) * ((R^(z+1)-1)/(R-1))
// NOD = P^x * Q^y * R^z  => here, P, Q, R are prime factors & x, y, z are
// powers NOD = (x + 1) (y + 1) (z + 1)
pair<int, int> SOD_NOD(int n) {
  int sod = 1, nod = 1;
  for (int i = 2; i * i <= n; ++i) {
    if (n % i == 0) {
      int pown = 1, pows = 0;
      while (n % i == 0) {
        pown *= i;  // p^e
        pows++;
        n /= i;
      }
      pown *= i;                    // p^e+1
      sod *= (pown - 1) / (i - 1);  //(p^e+1)-1 / p-1
      nod *= (pows + 1);
    }
  }
  if (n > 1) {
    sod *= (n + 1);
    nod *= 2;
  }
  return {sod, nod};
}
\end{minted}

\subsection{Segmented Sieve}
\begin{minted}[
    breaklines = true,
    breakanywhere = true,
    frame=lines,    
    fontfamily=tt,
    linenos=false,
    numberblanklines=true,
    numbersep=2pt,
    gobble=0,
    framerule=1pt,
    framesep=1mm,
    funcnamehighlighting=true,
    tabsize=2,
    obeytabs=false,
    mathescape=false
    samepage=false, %with this setting you can force the list to appear on the same page
    showspaces=false,
    showtabs =false,
    texcl=false
]{C++}

void segSeive(ll low, ll high) {
  vector < bool > area((high - low) + 1, true);
  for (ll i = 0; primes[i]*primes[i] <= high; i++) {
    ll start = ((low / primes[i]) * primes[i]);
    if (start < low) start += primes[i];
    for (ll j = start; j <= high; j += primes[i]) {
      if (j == primes[i]) continue;
      area[j - low] = false;
    }
  }
  for (ll i = 0; i < (high - low) + 1; i++) {
    if (area[i]) {
      if (i + low != 1 and i + low != 0) {
        cout << i + low << endl;
      }
    }
  }
}
\end{minted}

\subsection{Sieve}
\begin{minted}[
    breaklines = true,
    breakanywhere = true,
    frame=lines,    
    fontfamily=tt,
    linenos=false,
    numberblanklines=true,
    numbersep=2pt,
    gobble=0,
    framerule=1pt,
    framesep=1mm,
    funcnamehighlighting=true,
    tabsize=2,
    obeytabs=false,
    mathescape=false
    samepage=false, %with this setting you can force the list to appear on the same page
    showspaces=false,
    showtabs =false,
    texcl=false
]{C++}

const ll MAXN = 1e7 + 10;
bool prime[MAXN];
vector<ll> prm;
void sieve() {
  prime[0] = prime[1] = true;
  for (ll i = 2; i < MAXN; i++) {
    if (!prime[i]) {
      prm.push_back(i);
      for (ll j = i + i; j < MAXN; j += i) {
        prime[j] = true;
      }
    }
  }
}
\end{minted}

\subsection{Spf}
\begin{minted}[
    breaklines = true,
    breakanywhere = true,
    frame=lines,    
    fontfamily=tt,
    linenos=false,
    numberblanklines=true,
    numbersep=2pt,
    gobble=0,
    framerule=1pt,
    framesep=1mm,
    funcnamehighlighting=true,
    tabsize=2,
    obeytabs=false,
    mathescape=false
    samepage=false, %with this setting you can force the list to appear on the same page
    showspaces=false,
    showtabs =false,
    texcl=false
]{C++}

const int MAXN = 1e6 + 2;
int spf[MAXN];
vector<int> prms;
void preStore() {
  for (int i = 1; i < MAXN; i++) spf[i] = i;
  for (int i = 2; i < MAXN; i++) {
    if (spf[i] == i) {
      prms.push_back(i);
      for (int j = i + i; j < MAXN; j += i) {
        spf[j] = min(spf[j], i);
      }
    }
  }
}
\end{minted}

\subsection{UniquePF of all elements till MX}
\begin{minted}[
    breaklines = true,
    breakanywhere = true,
    frame=lines,    
    fontfamily=tt,
    linenos=false,
    numberblanklines=true,
    numbersep=2pt,
    gobble=0,
    framerule=1pt,
    framesep=1mm,
    funcnamehighlighting=true,
    tabsize=2,
    obeytabs=false,
    mathescape=false
    samepage=false, %with this setting you can force the list to appear on the same page
    showspaces=false,
    showtabs =false,
    texcl=false
]{C++}

const int MX = 2e5 + 10;
vector<int> pfac[MX];
void factorize() {
  for (int i = 2; i < MX; i++) {
    if (!pfac[i].empty()) continue;
    for (int j = i; j < MX; j += i)
      pfac[j].push_back(i);
  }
}
\end{minted}

\subsection{int128}
\begin{minted}[
    breaklines = true,
    breakanywhere = true,
    frame=lines,    
    fontfamily=tt,
    linenos=false,
    numberblanklines=true,
    numbersep=2pt,
    gobble=0,
    framerule=1pt,
    framesep=1mm,
    funcnamehighlighting=true,
    tabsize=2,
    obeytabs=false,
    mathescape=false
    samepage=false, %with this setting you can force the list to appear on the same page
    showspaces=false,
    showtabs =false,
    texcl=false
]{C++}

__int128 read() {
  __int128 x = 0, f = 1;
  char ch = getchar();
  while (ch < '0' || ch > '9') {
    if (ch == '-') f = -1;
    ch = getchar();
  }
  while (ch >= '0' && ch <= '9') {
    x = x * 10 + ch - '0';
    ch = getchar();
  }
  return x * f;
}
void print(__int128 x) {
  if (x < 0) {
    putchar('-');
    x = -x;
  }
  if (x > 9) print(x / 10);
  putchar(x % 10 + '0');
}
\end{minted}

\subsection{nCr and nPr}
\begin{minted}[
    breaklines = true,
    breakanywhere = true,
    frame=lines,    
    fontfamily=tt,
    linenos=false,
    numberblanklines=true,
    numbersep=2pt,
    gobble=0,
    framerule=1pt,
    framesep=1mm,
    funcnamehighlighting=true,
    tabsize=2,
    obeytabs=false,
    mathescape=false
    samepage=false, %with this setting you can force the list to appear on the same page
    showspaces=false,
    showtabs =false,
    texcl=false
]{C++}

int fact[N], ifact[N];
void prec() {
  fact[0] = 1;
  for (int i = 1; i < N; i++) {
    fact[i] = 1LL * fact[i - 1] * i % mod;
  }
  ifact[N - 1] = power(fact[N - 1], -1);
  for (int i = N - 2; i >= 0; i--) {
    ifact[i] = 1LL * ifact[i + 1] * (i + 1) % mod;
  }
}
int nPr(int n, int r) {
  if (n < r) return 0;
  return 1LL * fact[n] * ifact[n - r] % mod;
}
int nCr(int n, int r) {
  if (n < r) return 0;
  return 1LL * fact[n] * ifact[r] % mod * ifact[n - r] % mod;
}
\end{minted}

\subsection{nCr v2}
\begin{minted}[
    breaklines = true,
    breakanywhere = true,
    frame=lines,    
    fontfamily=tt,
    linenos=false,
    numberblanklines=true,
    numbersep=2pt,
    gobble=0,
    framerule=1pt,
    framesep=1mm,
    funcnamehighlighting=true,
    tabsize=2,
    obeytabs=false,
    mathescape=false
    samepage=false, %with this setting you can force the list to appear on the same page
    showspaces=false,
    showtabs =false,
    texcl=false
]{C++}

const int MX = 1e6 + 10;
const int M = 1e9 + 7;
int fact[MX], inv_fact[MX];
int modPow(int a, int b) {
  int ans = 1;
  while (b) {
    if (b & 1) ans = (1LL * ans * a) % M;
    a = (1LL * a * a) % M;
    b >>= 1;
  }
  return ans;
}
void precalFact() {
  fact[0] = inv_fact[0] = 1;
  for (int i = 1; i < MX; i++) {
    fact[i] = (1LL * fact[i - 1] * i) % M;
  }
  inv_fact[MX - 1] = modPow(fact[MX - 1], M - 2);
  for (int i = MX - 2; i >= 1; i--) {
    inv_fact[i] = (1LL * inv_fact[i + 1] * (i + 1)) % M;
  }
}
int nCr(int n, int r) {
  if (r < 0 or r > n) return 0;
  return 1LL * fact[n] * inv_fact[r] % M * inv_fact[n - r] % M;
}
\end{minted}

\section{String}
\subsection{Aho Corasic}
\begin{minted}[
    breaklines = true,
    breakanywhere = true,
    frame=lines,    
    fontfamily=tt,
    linenos=false,
    numberblanklines=true,
    numbersep=2pt,
    gobble=0,
    framerule=1pt,
    framesep=1mm,
    funcnamehighlighting=true,
    tabsize=2,
    obeytabs=false,
    mathescape=false
    samepage=false, %with this setting you can force the list to appear on the same page
    showspaces=false,
    showtabs =false,
    texcl=false
]{C++}

//number of occourence of word in a text
const ll N = 1e6+10, A = 26;
ll trie[N][A], pos[N], slink[N], dp[N], tot = 1;
vector<int> order;
void initTrie(){
  order.clear();
  while(tot--){
    memset(trie[tot], 0, sizeof(trie[tot]));
  }
  memset(pos, 0, sizeof(pos));
  memset(slink, 0, sizeof(slink));
  memset(dp, 0, sizeof(dp)); tot = 1;
}
void addStr(string &s, int ind){
  ll u = 0;
  for(auto it: s){
    ll n = it-'a';
    if(trie[u][n]==0) trie[u][n] = tot++;
    u = trie[u][n];
  } pos[ind] = u;
}
void build(){
  queue<ll> q; q.push(0);
  while(!q.empty()){
    ll p = q.front(); q.pop();
    order.push_back(p);
    for(ll c = 0; c<A; c++){
      ll u = trie[p][c];
      if(!u) continue;
      q.push(u);
      if(!p) continue;
      ll v = slink[p];
      while(v && !trie[v][c]) v = slink[v];
      slink[u] = trie[v][c];
    }
  }
}
void trav(string &s){
  ll u = 0;
  for(char c: s){
    c-='a';
    while(u && !trie[u][c]) u = slink[u];
    u = trie[u][c]; dp[u]++;
  }
  reverse(order.begin(), order.end());
  for(auto u: order){
    dp[slink[u]]+=dp[u];
  }
}
void solve(){
  ll n; cin>>n;
  string text; cin>>text;
  string s;
  for(ll i = 0; i<n; i++){
    cin>>s; addStr(s, i);
  }
  build(); trav(text);
  for(ll i = 0; i<n; i++){
    cout<<dp[pos[i]]<<endl;
  }
}
int32_t main(){
  ios_base::sync_with_stdio(0);
  cin.tie(0);
  ll tc = 1;
  cin>>tc;
  for(ll i = 1; i<=tc;i++){
    cout<<"Case "<<i<<": \n";
    initTrie();
    solve();
  }
}
\end{minted}

\subsection{LCS for 3 Strings}
\begin{minted}[
    breaklines = true,
    breakanywhere = true,
    frame=lines,    
    fontfamily=tt,
    linenos=false,
    numberblanklines=true,
    numbersep=2pt,
    gobble=0,
    framerule=1pt,
    framesep=1mm,
    funcnamehighlighting=true,
    tabsize=2,
    obeytabs=false,
    mathescape=false
    samepage=false, %with this setting you can force the list to appear on the same page
    showspaces=false,
    showtabs =false,
    texcl=false
]{C++}

string a, b, c;
ll dp[55][55][55];
ll lcs(ll i, ll j, ll k) {
  if (i == a.size() or j == b.size() or k == c.size()) return 0;
  if (dp[i][j][k] != -1) return dp[i][j][k];
  if (a[i] == b[j] and a[i] == c[k]) return 1 + lcs(i + 1, j + 1, k + 1);
  ll ans = 0;
  ans = max(ans, lcs(i, j , k + 1));
  ans = max(ans, lcs(i, j + 1 , k));
  ans = max(ans, lcs(i + 1, j, k));
  return dp[i][j][k] = ans;
}
\end{minted}

\subsection{Manacher Palindrome}
\begin{minted}[
    breaklines = true,
    breakanywhere = true,
    frame=lines,    
    fontfamily=tt,
    linenos=false,
    numberblanklines=true,
    numbersep=2pt,
    gobble=0,
    framerule=1pt,
    framesep=1mm,
    funcnamehighlighting=true,
    tabsize=2,
    obeytabs=false,
    mathescape=false
    samepage=false, %with this setting you can force the list to appear on the same page
    showspaces=false,
    showtabs =false,
    texcl=false
]{C++}

// pal[1][i] = longest odd (half rounded down) palindrome around pos i and
// starts at i - pal[1][i] and ends at i + pal[1][i] pal[0][i] = half length of
// longest even palindrome around pos i, i + 1 and starts at i - par[0][i] + 1
// and ends at i + pal[0][i]
const int N = 5e5 + 10;
int pal[2][N];
void manacher(string& s) {
  int n = s.size(), idx = 2;
  while (idx--) {
    for (int l = -1, r = -1, i = 0; i < n - 1; ++i) {
      if (i > r)
        l = r = i;
      else {
        int k = min(r - i, pal[idx][l + r - i]);
        l = i - k, r = i + k;
      }
      while (l - idx >= 0 and r + 1 < n and s[l - idx] == s[r + 1]) l--, r++;
      pal[idx][i] = r - i;
      // [l - 1 + idx : r] palindrome
    }
  }
}
\end{minted}

\subsection{String Hashing 2}
\begin{minted}[
    breaklines = true,
    breakanywhere = true,
    frame=lines,    
    fontfamily=tt,
    linenos=false,
    numberblanklines=true,
    numbersep=2pt,
    gobble=0,
    framerule=1pt,
    framesep=1mm,
    funcnamehighlighting=true,
    tabsize=2,
    obeytabs=false,
    mathescape=false
    samepage=false, %with this setting you can force the list to appear on the same page
    showspaces=false,
    showtabs =false,
    texcl=false
]{C++}

const int N = 1000010, MOD = 1e9 + 7;
const ll P[] = {97, 1000003};
ll bigMod(ll a, ll e) {
  if (e == -1) e = MOD - 2;
  ll ret = 1;
  while (e) {
    if (e & 1) ret = ret * a % MOD;
    a = a * a % MOD, e >>= 1;
  }
  return ret;
}
ll pwr[2][N], inv[2][N];
void initHash() {
  for (int it = 0; it < 2; ++it) {
    pwr[it][0] = inv[it][0] = 1;
    ll INV_P = bigMod(P[it], -1);
    for (int i = 1; i < N; ++i) {
      pwr[it][i] = pwr[it][i - 1] * P[it] % MOD;
      inv[it][i] = inv[it][i - 1] * INV_P % MOD;
    }
  }
}
struct RangeHash {
  vector<ll> h[2], rev[2];
  RangeHash(const string S, bool revFlag = 0) {
    for (int it = 0; it < 2; ++it) {
      h[it].resize(S.size() + 1, 0);
      for (int i = 0; i < S.size(); ++i) {
        h[it][i + 1] = (h[it][i] + pwr[it][i + 1] * (S[i] - 'a' + 1)) % MOD;
      }
      if (revFlag) {
        rev[it].resize(S.size() + 1, 0);
        for (int i = 0; i < S.size(); ++i) {
          rev[it][i + 1] =
              (rev[it][i] + inv[it][i + 1] * (S[i] - 'a' + 1)) % MOD;
        }
      }
    }
  }
  inline ll get(int l, int r) {
    ll one = (h[0][r + 1] - h[0][l]) * inv[0][l + 1] % MOD;
    ll two = (h[1][r + 1] - h[1][l]) * inv[1][l + 1] % MOD;
    if (one < 0) one += MOD;
    if (two < 0) two += MOD;
    return one << 31 | two;
  }
  inline ll getReverse(int l, int r) {
    ll one = (rev[0][r + 1] - rev[0][l]) * pwr[0][r + 1] % MOD;
    ll two = (rev[1][r + 1] - rev[1][l]) * pwr[1][r + 1] % MOD;
    if (one < 0) one += MOD;
    if (two < 0) two += MOD;
    return one << 31 | two;
  }
};
\end{minted}

\subsection{String Hashing}
\begin{minted}[
    breaklines = true,
    breakanywhere = true,
    frame=lines,    
    fontfamily=tt,
    linenos=false,
    numberblanklines=true,
    numbersep=2pt,
    gobble=0,
    framerule=1pt,
    framesep=1mm,
    funcnamehighlighting=true,
    tabsize=2,
    obeytabs=false,
    mathescape=false
    samepage=false, %with this setting you can force the list to appear on the same page
    showspaces=false,
    showtabs =false,
    texcl=false
]{C++}

const int mod1 = 911382323, mod2 = 972663749, b1 = 137, b2 = 139;
const int mxN = 5000010;
int pow_b1[mxN], pow_b2[mxN], inv_b1[mxN], inv_b2[mxN];
int binExp(int base, int power, int mod) {
  int res = 1;
  while (power) {
    if (power & 1) res = (1LL * res * base) % mod;
    base = (1LL * base * base) % mod;
    power >>= 1;
  }
  return res;
}
void pre() {
  pow_b1[0] = pow_b2[0] = 1;
  for (int i = 1; i < mxN; i++) {
    pow_b1[i] = (1LL * pow_b1[i - 1] * b1) % mod1;
    pow_b2[i] = (1LL * pow_b2[i - 1] * b2) % mod2;
  }
  inv_b1[mxN - 1] = binExp(pow_b1[mxN - 1], mod1 - 2, mod1);
  inv_b2[mxN - 1] = binExp(pow_b2[mxN - 1], mod2 - 2, mod2);
  for (int i = mxN - 2; i >= 0; i--) {
    inv_b1[i] = (1LL * inv_b1[i + 1] * b1) % mod1;
    inv_b2[i] = (1LL * inv_b2[i + 1] * b2) % mod2;
  }
}
vector<pair<int, int>> getPref(string& s) {
  int qq = s.size();
  vector<pair<int, int>> hsh(qq);
  for (int i = 0; i < qq; i++) {
    if (i == 0) {
      hsh[i].first = (1LL * s[i] * pow_b1[i]) % mod1;
      hsh[i].second = (1LL * s[i] * pow_b2[i]) % mod2;
    } else {
      hsh[i].first =
          (hsh[i - 1].first + (1LL * s[i] * pow_b1[i]) % mod1) % mod1;
      hsh[i].second =
          (hsh[i - 1].second + (1LL * s[i] * pow_b2[i]) % mod2) % mod2;
    }
  }
  return hsh;
}
pair<int, int> getHash(string& str) {
  int hsh1 = 0, hsh2 = 0, sz = str.size();
  for (int i = 0; i < sz; ++i) {
    hsh1 = (hsh1 + 1LL * str[i] * pow_b1[i] % mod1) % mod1;
  }
  for (int i = 0; i < sz; ++i) {
    hsh2 = (hsh2 + 1LL * str[i] * pow_b2[i] % mod2) % mod2;
  }
  return {hsh1, hsh2};
}
pair<int, int> getSub(int l, int r, vector<pair<int, int>>& v) {
  pair<int, int> q;
  if (l == 0) {
    q = {v[r].first, v[r].second};
  } else {
    int x = (1LL * ((v[r].first - v[l - 1].first + mod1) % mod1) * inv_b1[l]) %
            mod1;
    int y =
        (1LL * ((v[r].second - v[l - 1].second + mod2) % mod2) * inv_b2[l]) %
        mod2;
    q = {x, y};
  }
  return q;
}
\end{minted}

\subsection{Suffix Array}
\begin{minted}[
    breaklines = true,
    breakanywhere = true,
    frame=lines,    
    fontfamily=tt,
    linenos=false,
    numberblanklines=true,
    numbersep=2pt,
    gobble=0,
    framerule=1pt,
    framesep=1mm,
    funcnamehighlighting=true,
    tabsize=2,
    obeytabs=false,
    mathescape=false
    samepage=false, %with this setting you can force the list to appear on the same page
    showspaces=false,
    showtabs =false,
    texcl=false
]{C++}

// fahimcp495
array<vector<int>, 2> get_sa(string& s, int lim = 128) {  // for integer, just change string to vector<int> and minimum value of vector must be >= 1
  int n = s.size() + 1, k = 0, a, b;
  vector<int> x(begin(s), end(s) + 1), y(n), sa(n), lcp(n), ws(max(n, lim)), rank(n);
  x.back() = 0;
  iota(begin(sa), end(sa), 0);
  for (int j = 0, p = 0; p < n; j = max(1, j * 2), lim = p) {
    p = j, iota(begin(y), end(y), n - j);
    for (int i = 0; i < n; ++i)
      if (sa[i] >= j) y[p++] = sa[i] - j;
    fill(begin(ws), end(ws), 0);
    for (int i = 0; i < n; ++i) ws[x[i]]++;
    for (int i = 1; i < lim; ++i) ws[i] += ws[i - 1];
    for (int i = n; i--;) sa[--ws[x[y[i]]]] = y[i];
    swap(x, y), p = 1, x[sa[0]] = 0;
    for (int i = 1; i < n; ++i) a = sa[i - 1], b = sa[i], x[b] = (y[a] == y[b] && y[a + j] == y[b + j]) ? p - 1 : p++;
  }
  for (int i = 1; i < n; ++i) rank[sa[i]] = i;
  for (int i = 0, j; i < n - 1; lcp[rank[i++]] = k)
    for (k &&k--, j = sa[rank[i] - 1]; s[i + k] == s[j + k]; k++);
  sa.erase(sa.begin()), lcp.erase(lcp.begin());
  return {sa, lcp};
}
\end{minted}

\subsection{Suffix Automata}
\begin{minted}[
    breaklines = true,
    breakanywhere = true,
    frame=lines,    
    fontfamily=tt,
    linenos=false,
    numberblanklines=true,
    numbersep=2pt,
    gobble=0,
    framerule=1pt,
    framesep=1mm,
    funcnamehighlighting=true,
    tabsize=2,
    obeytabs=false,
    mathescape=false
    samepage=false, %with this setting you can force the list to appear on the same page
    showspaces=false,
    showtabs =false,
    texcl=false
]{C++}

const int N = 2e5 + 10;  // max string size
int len[N], lnk[N]{-1}, last, sz = 1;
unordered_map<char, int> to[N];
void add(char c) {
  int cur = sz++;
  len[cur] = len[last] + 1;
  int u = last;
  while (u != -1 and !to[u].count(c)) {
    to[u][c] = cur;
    u = lnk[u];
  }
  if (u == -1) {
    lnk[cur] = 0;
  } else {
    int v = to[u][c];
    if (len[v] == len[u] + 1) {
      lnk[cur] = v;
    } else {
      int w = sz++;
      len[w] = len[u] + 1, lnk[w] = lnk[v], to[w] = to[v];
      while (u != -1 and to[u][c] == v) {
        to[u][c] = w;
        u = lnk[u];
      }
      lnk[cur] = lnk[v] = w;
    }
  }
  last = cur;
}
\end{minted}

\subsection{Suffix Automation}
\begin{minted}[
    breaklines = true,
    breakanywhere = true,
    frame=lines,    
    fontfamily=tt,
    linenos=false,
    numberblanklines=true,
    numbersep=2pt,
    gobble=0,
    framerule=1pt,
    framesep=1mm,
    funcnamehighlighting=true,
    tabsize=2,
    obeytabs=false,
    mathescape=false
    samepage=false, %with this setting you can force the list to appear on the same page
    showspaces=false,
    showtabs =false,
    texcl=false
]{C++}

int len[N], lnk[N]{-1}, last, sz = 1;
unordered_map<char, int> to[N];
void init() {
  while (sz) {
    sz--;
    to[sz].clear();
  }
  last = 0, sz = 1;
}
void add(char c) {
  int cur = sz++;
  int u = last;
  len[cur] = len[last] + 1;
  while (u != -1 and !to[u].count(c)) {
    to[u][c] = cur;
    u = lnk[u];
  }
  if (u == -1) {
    lnk[cur] = 0;
  } else {
    int v = to[u][c];
    if (len[v] == len[u] + 1) {
      lnk[cur] = v;
    } else {
      int w = sz++;
      len[w] = len[u] + 1, lnk[w] = lnk[v], to[w] = to[v];
      while (u != -1 and to[u][c] == v) {
        to[u][c] = w;
        u = lnk[u];
      }
      lnk[cur] = lnk[v] = w;
    }
  }
  last = cur;
}
\end{minted}

\subsection{Trie}
\begin{minted}[
    breaklines = true,
    breakanywhere = true,
    frame=lines,    
    fontfamily=tt,
    linenos=false,
    numberblanklines=true,
    numbersep=2pt,
    gobble=0,
    framerule=1pt,
    framesep=1mm,
    funcnamehighlighting=true,
    tabsize=2,
    obeytabs=false,
    mathescape=false
    samepage=false, %with this setting you can force the list to appear on the same page
    showspaces=false,
    showtabs =false,
    texcl=false
]{C++}

const ll N = 1e6 + 5, A = 26;
ll trie[N][A], cnt[N], tot = 1, root = 1;
void initTrie() {
  cnt[tot] = 0;
  root = 1;
}
void addStr(string& s) {
  ll u = 1;
  for (auto it : s) {
    ll n = it - 'a';
    if (trie[u][n] == 0) {
      trie[u][n] = ++tot;
    }
    u = trie[u][n];
    cnt[u]++;
  }
}
ll wordCount(string& s) {
  ll u = 1;
  for (auto it : s) {
    int n = it - 'a';
    if (trie[u][n] == 0) return 0;
    u = trie[u][n];
  }
  return cnt[u];
}
\end{minted}

\section{Tree}
\subsection{Centroid Decomposition}
\begin{minted}[
    breaklines = true,
    breakanywhere = true,
    frame=lines,    
    fontfamily=tt,
    linenos=false,
    numberblanklines=true,
    numbersep=2pt,
    gobble=0,
    framerule=1pt,
    framesep=1mm,
    funcnamehighlighting=true,
    tabsize=2,
    obeytabs=false,
    mathescape=false
    samepage=false, %with this setting you can force the list to appear on the same page
    showspaces=false,
    showtabs =false,
    texcl=false
]{C++}

const int N = 2e5+5;
int n, k, sz[N], centered[N], ans = 0;
vector<int> adj[N];
void dfs_sz(int u, int p) {
  sz[u] = 1;
  for (auto v: adj[u]) {
    if (v != p && !centered[v]) {
      dfs_sz(v, u); sz[u] += sz[v];
    }
  }
}
int get_cen(int u, int p, int n) {
  for (auto v: adj[u]) {
    if (v != p && !centered[v] && sz[v] > n/2) {
      return get_cen(v, u, n);
    }
  } return u;
}
int t, tin[N], tout[N], nodes[N], dis[N];
void dfs(int u, int p){
  nodes[t] = u;
  tin[u] = t++;
  for(auto v: adj[u]){
    if(v!=p && !centered[v]){
      dis[v] = dis[u]+1; dfs(v, u);
    }
  } tout[u] = t-1;
}
void go(int u){
  dfs_sz(u, u);
  int c = get_cen(u, u, sz[u]);
  centered[c] = 1; sz[c] = sz[u];
  t = 0; dis[c] = 0; dfs(c, c);
  int cnt[t]{1};
  for(auto v: adj[c]){
    if(centered[v]) continue;
    for(int i = tin[v]; i<=tout[v]; ++i){
      int w = nodes[i];
      if(k-dis[w]>=0 && k-dis[w]<t){
        ans+=cnt[k-dis[w]];
      }
    }
    for(int i = tin[v]; i<=tout[v]; ++i){
      int w = nodes[i]; cnt[dis[w]]++;
    }
  }
  for(auto v: adj[c]){
    if(!centered[v]){  go(v); }
  }
}
void solve() {
  cin>>n>>k;
  for(ll i = 1; i<n; i++){
    ll u, v; cin>>u>>v;
    adj[u].push_back(v);
    adj[v].push_back(u);
  } go(1);
  cout<<ans<<endl;
}
\end{minted}

\subsection{DSUOnTrees}
\begin{minted}[
    breaklines = true,
    breakanywhere = true,
    frame=lines,    
    fontfamily=tt,
    linenos=false,
    numberblanklines=true,
    numbersep=2pt,
    gobble=0,
    framerule=1pt,
    framesep=1mm,
    funcnamehighlighting=true,
    tabsize=2,
    obeytabs=false,
    mathescape=false
    samepage=false, %with this setting you can force the list to appear on the same page
    showspaces=false,
    showtabs =false,
    texcl=false
]{C++}

int n, color[MX], ans[MX];
vector<int> g[MX];
set<int> bucket[MX];
int merge(int a, int b) {
  if (bucket[a].size() < bucket[b].size()) swap(a, b);
  bucket[a].insert(bucket[b].begin(), bucket[b].end());
  bucket[b].clear();
  return a;
}
int dfs(int u, int p = -1) {
  int cur = u;
  for (int v : g[u])
    if (v != p)
      cur = merge(cur, dfs(v, u));
  ans[u] = (int)bucket[cur].size();
  return cur;
}
void solve() {
  cin >> n;
  for (int i = 0; i < n; ++i) {
    cin >> color[i];
    bucket[i].insert(color[i]);
  }
  // graph input
  dfs(0);
  // print output
}
\end{minted}

\subsection{LCA using binary Lifting}
\begin{minted}[
    breaklines = true,
    breakanywhere = true,
    frame=lines,    
    fontfamily=tt,
    linenos=false,
    numberblanklines=true,
    numbersep=2pt,
    gobble=0,
    framerule=1pt,
    framesep=1mm,
    funcnamehighlighting=true,
    tabsize=2,
    obeytabs=false,
    mathescape=false
    samepage=false, %with this setting you can force the list to appear on the same page
    showspaces=false,
    showtabs =false,
    texcl=false
]{C++}

int n, l;
vector<vector<int>> adj;
int timer;
vector<int> tin, tout;
vector<vector<int>> up;
void dfs(int v, int p) {
  tin[v] = ++timer;
  up[v][0] = p;
  for (int i = 1; i <= l; ++i)
    up[v][i] = up[up[v][i - 1]][i - 1];

  for (int u : adj[v]) {
    if (u != p)
      dfs(u, v);
  }
  tout[v] = ++timer;
}
bool is_ancestor(int u, int v) {
  return tin[u] <= tin[v] && tout[u] >= tout[v];
}
int lca(int u, int v) {
  if (is_ancestor(u, v))
    return u;
  if (is_ancestor(v, u))
    return v;
  for (int i = l; i >= 0; --i) {
    if (!is_ancestor(up[u][i], v))
      u = up[u][i];
  }
  return up[u][0];
}
void preprocess(int root) {
  tin.resize(n);
  tout.resize(n);
  timer = 0;
  l = ceil(log2(n));
  up.assign(n, vector<int>(l + 1));
  dfs(root, root);
}
\end{minted}

\subsection{LCA}
\begin{minted}[
    breaklines = true,
    breakanywhere = true,
    frame=lines,    
    fontfamily=tt,
    linenos=false,
    numberblanklines=true,
    numbersep=2pt,
    gobble=0,
    framerule=1pt,
    framesep=1mm,
    funcnamehighlighting=true,
    tabsize=2,
    obeytabs=false,
    mathescape=false
    samepage=false, %with this setting you can force the list to appear on the same page
    showspaces=false,
    showtabs =false,
    texcl=false
]{C++}

const int N = 1e5 + 5;
vector<int> g[N], parent(N), depth(N, 0);
void dfs(int vertex, int par = -1) {
  parent[vertex] = par;
  for (auto child : g[vertex]) {
    if (child != par) {
        depth[child] = depth[vertex] + 1;  dfs(child, vertex);
    }
  }
}
int lca(int x, int y) {
  int diff = min(depth[x], depth[y]);
  while (depth[x] > diff) x = parent[x];
  while (depth[y] > diff) y = parent[y];
  while (x != y) { x = parent[x]; y = parent[y]; }  return x;
}
\end{minted}

\subsection{Tree Rerooting(A simple Path)}
\begin{minted}[
    breaklines = true,
    breakanywhere = true,
    frame=lines,    
    fontfamily=tt,
    linenos=false,
    numberblanklines=true,
    numbersep=2pt,
    gobble=0,
    framerule=1pt,
    framesep=1mm,
    funcnamehighlighting=true,
    tabsize=2,
    obeytabs=false,
    mathescape=false
    samepage=false, %with this setting you can force the list to appear on the same page
    showspaces=false,
    showtabs =false,
    texcl=false
]{C++}

#include <bits/stdc++.h>
#define int long long
#define endl '\n'
using namespace std;
const int N = 3e5+5;
vector < int > tree[N];
int ar[N];
int dp1[N];
int ans[N];
int n,k,q;
void clear(int n) {
    for(int i = 0; i <= n; i++) {
      tree[i].clear();
      ar[i] = 0;
      ans[i] = 0;
      dp1[i] = 0;
    }
}
void pre(int vertex, int par) {
    dp1[vertex] = 0;
    for(auto child : tree[vertex]) {
      if(child == par) continue;
      pre(child,vertex);
      dp1[vertex] = max(dp1[vertex],dp1[child]);

    }
    dp1[vertex] += ar[vertex];
}
void reroot(int vertex, int par, int max_val) {
    ans[vertex] = max(dp1[vertex],max_val+ar[vertex]);
    int maxi = 0, low = 0;
    for(auto child : tree[vertex]) {
      if(child == par) continue;
      int val = dp1[child];
      if(val > maxi) {
            low = maxi;
            maxi = val;
      }
      else if(val > low) {
            low = val;
      }
    }
    for(auto child : tree[vertex]) {
      if(child == par) continue;
      int x = (dp1[child] == maxi) ? low : maxi;
      int boro = max(max_val,x) + ar[vertex];

      reroot(child,vertex,boro);
    }
}
void solve()
{
    cin >> n >> k >> q;
    clear(n);
    for(int i = 1; i <= k; i++) {
      int x; cin >> x;
      ar[x] = 1;
    }
    for(int i = 1; i < n; i++) {
      int u,v; cin >> u >> v;
      tree[u].push_back(v);
      tree[v].push_back(u);
    }
    vector < int > qr;
    while(q--) {
      int x; cin >> x;
      qr.push_back(x);
    }
    pre(1,-1);
    reroot(1,-1,0);
    int maxi = 0;
    for(int i = 1; i <= n; i++) {
      maxi = max(maxi,ans[i]);
    }
    for(auto &it:qr) {
      if(ans[it] == maxi) cout << "JA" << endl;
      else cout << "NEIN" << endl;
    }
}
signed main()
{
    ios_base::sync_with_stdio(0),cin.tie(0);
    int tt = 1;
    cin >> tt;
    while(tt--)
    solve();
}
\end{minted}

\subsection{Tree Rerooting}
\begin{minted}[
    breaklines = true,
    breakanywhere = true,
    frame=lines,    
    fontfamily=tt,
    linenos=false,
    numberblanklines=true,
    numbersep=2pt,
    gobble=0,
    framerule=1pt,
    framesep=1mm,
    funcnamehighlighting=true,
    tabsize=2,
    obeytabs=false,
    mathescape=false
    samepage=false, %with this setting you can force the list to appear on the same page
    showspaces=false,
    showtabs =false,
    texcl=false
]{C++}

#include <bits/stdc++.h>
using namespace std;

const int N = 2e5+5;
vector < int > tree[N];
vector < int > subtree_sum(N,1);
vector < int > dist(N,0);
int ans[N];
int ar[N];
int n;

void pre(int vertex, int par) {

    subtree_sum[vertex] = ar[vertex];
    for(auto child : tree[vertex]) {
      if(child == par) continue;
      pre(child,vertex);
      dist[vertex] += dist[child] + subtree_sum[child];
      subtree_sum[vertex] += subtree_sum[child];
    }
}

void reroot(int vertex, int par) {
    ans[vertex] = dist[vertex];
    for(auto child : tree[vertex]) {
      if(child == par) continue;

      dist[vertex] -= (dist[child] + subtree_sum[child]);
      subtree_sum[vertex] -= subtree_sum[child];

      dist[child] += (dist[vertex] + subtree_sum[vertex]);
      subtree_sum[child] += subtree_sum[vertex];

      reroot(child,vertex);

      dist[child] -= (dist[vertex] + subtree_sum[vertex]);
      subtree_sum[child] -= subtree_sum[vertex];

      dist[vertex] += (dist[child] + subtree_sum[child]);
      subtree_sum[vertex] += subtree_sum[child];
    }
}
void solve()
{
    cin >> n;
    for(int i = 1; i <= n; i++) {
      cin >> ar[i];
    }
    for(int i = 1; i < n; i++) {
      int u,v; cin >> u >> v;
      tree[u].push_back(v);
      tree[v].push_back(u);
    }

    pre(1,-1);
    reroot(1,-1);
    int maxi = 0;
    for(int i = 1; i <= n; i++) {
      maxi = max(maxi,ans[i]);
    }
    cout << maxi << endl;
}
signed main()
{
    int tt = 1;
    while(tt--)
    solve();
}
\end{minted}

\section{Notes}
\subsection{Geometry}

\subsubsection{Triangles}
Circumradius: $R=\dfrac{abc}{4A}$ , Inradius: $r=\dfrac{A}{s}$\\
The area of a triangle using two sides and the included angle can be given as:  \[ A = \frac{1}{2} ab \sin \angle C \]
Length of median (divides triangle into two equal-area triangles): $m_a=\tfrac{1}{2}\sqrt{2b^2+2c^2-a^2}$\\
Length of bisector (divides angles in two): $s_a=\sqrt{bc\left[1-\left(\dfrac{a}{b+c}\right)^2\right]}$\\
Law of tangents: $\dfrac{a+b}{a-b}=\dfrac{\tan\dfrac{\alpha+\beta}{2}}{\tan\dfrac{\alpha-\beta}{2}}$\\
Area given 3 sides of a triangle $a, b, c$: $A=\sqrt{s(s-a)(s-b)(s-c)}$ where $s=\dfrac{a+b+c}{2}$\\

\subsubsection{Sphere}
Volume of a sphere: $V = \frac{4}{3}\pi R^3$ \\
Spherical Cap volume: $V(h) = \frac{\pi h^2 (3R - h)}{3}$

\subsubsection{Quadrilaterals}
With side lengths $a,b,c,d$, diagonals $e, f$, diagonals angle $\theta$, area $A$ and
magic flux $F=b^2+d^2-a^2-c^2$:
\[ 4A = 2ef \cdot \sin\theta = F\tan\theta = \sqrt{4e^2f^2-F^2} \]

For cyclic quadrilaterals the sum of opposite angles is $180^\circ$,
$ef = ac + bd$, and $A = \sqrt{(p-a)(p-b)(p-c)(p-d)}$.

\subsubsection{Spherical coordinates}
\centerline{\includegraphics[width=25mm]{../code/sphericalCoordinates}}
\[
  \begin{array}{cc}
    x = r\sin\theta\cos\phi & r = \sqrt{x^2+y^2+z^2}\\
    y = r\sin\theta\sin\phi & \theta = \textrm{acos}(z/\sqrt{x^2+y^2+z^2})\\
    z = r\cos\theta & \phi = \textrm{atan2}(y,x)
\end{array}\]

\subsubsection{Pick's Theorem}
Given a lattice polygon with non-zero area, we define: $S$ as the area of the polygon, $I$ as the number of integer-coordinate points strictly inside the polygon, $B$ as the number of integer-coordinate points on the boundary of the polygon.
Then, Pick's Theorem states: $$S=I+\frac{B}{2}-1$$
The number of lattice points on segments $(x_1, y_1)$ to $(x_2 , y_2)$ is: $\gcd(|x_2 - x_1| , |y_2 - y_1|) + 1$

\subsubsection{Polygon}
For a regular polygon with \( n \) sides and side length \( a \), the circumradius \( R \) is given by:
\[
  R = \frac{a}{2 \sin\left(\frac{\pi}{n}\right)}
\]

\subsubsection{Area of a Circular Segment}
The area of a circular segment, which is the region enclosed by a chord and the corresponding arc, can be calculated using the formula:
\[
  A = \frac{R^2}{2} \left(\theta - \sin\theta\right)
\]
where:
\( R \) is the radius of the circle,
\( \theta \) is the central angle subtended by the chord, in radians.

\subsubsection{Matrix Rotation}
Anti-Clockwise Rotation:
$$
\begin{bmatrix}
\cos \theta & \sin \theta \\
-\sin \theta & \cos \theta
\end{bmatrix}
\begin{bmatrix}
x \\
y
\end{bmatrix}
=
\begin{bmatrix}
x' \\
y'
\end{bmatrix}
$$

Clockwise Rotation:
$$
\begin{bmatrix}
\cos \theta & -\sin \theta \\
\sin \theta & \cos \theta
\end{bmatrix}
$$

\subsection{Number Theory}

\subsubsection{Primes}
$p=962592769$ is such that $2^{21} \mid p-1$, which may be useful. For hashing
use 970592641 (31-bit number), 31443539979727 (45-bit), 3006703054056749
(52-bit). There are 78498 primes less than 1\,000\,000.

Primitive roots exist modulo any prime power $p^a$, except for $p = 2, a > 2$, and there are $\phi(\phi(p^a))$ many.
For $p = 2, a > 2$, the group $\mathbb Z_{2^a}^\times$ is instead isomorphic to $\mathbb Z_2 \times \mathbb Z_{2^{a-2}}$.

\subsubsection{Estimates}
$\sum_{d|n} d = O(n \log \log n)$.

The number of divisors of $n$ is at most around 100 for $n < 5e4$, 500 for $n < 1e7$, 2000 for $n < 1e10$, 200\,000 for $n < 1e19$.

\subsubsection{Perfect numbers}
$n>1$ is called perfect if it equals
sum of its proper divisors and $1$.  Even $n$ is perfect iff $n = 2^{p-1} (2^p - 1)$
and $2^p - 1$ is prime (Mersenne's). No odd perfect numbers are yet found.

\subsubsection{Carmichael numbers}
A positive composite $n$ is a Carmichael number
($a^{n-1} \equiv 1 \pmod{n}$ for all $a$: $\gcd(a,n)=1$),
iff $n$ is square-free, and for all prime divisors $p$ of $n$, $p-1$ divides $n-1$.

\subsubsection{Totient}
\begin{itemize}
  \item If \( p \) is a prime \( \phi(p^k) = p^k - p^{k-1} \)
  \item If \( a \) \& \( b \) are relatively prime, \( \phi(ab) = \phi(a)\phi(b) \)
  \item \( \phi(n) = n(1-\frac{1}{p_1})(1-\frac{1}{p_2})(1-\frac{1}{p_3})\dots(1-\frac{1}{p_k}) \)
  \item Sum of coprime to \( n = n \times \frac{\phi(n)}{2} \)
  \item If \( n = 2^k, \phi(n) = 2^{k - 1} = \frac{n}{2} \)
  \item For \( a \) \& \( b \), \( \phi(ab) = \phi(a)\phi(b)\frac{d}{\phi(d)} \)
  \item \( \phi (ip) = p \phi(i) \) whenever \( p \) is a prime and it divides \( i \)
  \item The number of \( a (1 \le a \le N) \) such that \( \gcd(a, N)=d \) is \( \phi(\frac{n}{d}) \)
  \item If \( n > 2 \) , \( \phi(n) \) is always even
  \item Sum of gcd, \( \sum_{i=1}^n \gcd(i, n) = \sum_{d|n} d \phi(\frac{n}{d}) \)
  \item Sum of lcm, \( \sum_{i=1}^{n}\operatorname{lcm}(i, n) = \frac{n}{2}(\sum_{d|n}(d \phi(d))+1) \)
  \item \( \phi(1) = 1 \) and \( \phi(2) = 1 \) which two are only odd \( \phi \)
  \item \( \phi(3) = 2 \) and \( \phi(4) = 2 \) and \( \phi(6) = 2 \) which three are only prime \( \phi \)
  \item Find minimum $n$ such that $ \frac{\phi(n)} {n} $ is  maximum - Multiple of small primes - $ 2 \cdot 3 \cdot 5 \cdot 7 \cdot 11 \cdot 13 \dots $
\end{itemize}

\subsubsection{Mobius function}
$\mu(1) = 1$. $\mu(n) = 0$, if $n$ is not squarefree.
$\mu(n) = (-1)^s$, if $n$ is the product of $s$ distinct primes.
Let $f$, $F$ be functions on positive integers.
If for all $n \in \mathbb{N}$, $F(n)=\sum_{d|n} f(d)$, then $f(n) = \sum_{d|n} \mu(d) F(\frac{n}{d})$,
and vice versa. \quad
$\phi(n) = \sum_{d|n} \mu(d) \frac{n}{d}$.
\quad $\sum_{d|n} \mu(d) = 1$. \\
If $f$ is multiplicative, then $\sum_{d|n} \mu(d) f(d) = \prod_{p|n}(1-f(p))$,
$\sum_{d|n} \mu(d)^2 f(d) = \prod_{p|n} (1+f(p))$.

\[ \sum_{i = 1}^n \sum_{j = 1}^n [\gcd(i, j) = 1] = \sum_{k = 1}^n \mu(k) \lfloor \frac{n}{k} \rfloor^2 \]
\[ \sum_{i = 1}^n \sum_{j = 1}^n \gcd(i, j) = \sum_{k = 1}^n k \sum_{l = 1}^{\lfloor \frac{n}{k} \rfloor} \mu(l) \lfloor {\frac{n}{kl}} \rfloor^2 \]
\[ \sum_{i = 1}^n \sum_{j = 1}^n \gcd(i, j) = \sum_{k = 1}^n \left(\frac{\lfloor \frac{n}{k} \rfloor (1 + \lfloor \frac{n}{k} \rfloor)}{2}\right)^2 \sum_{d | k} \mu(d) \frac{k}{d} \]

\subsubsection{Legendre symbol}
If $p$ is an odd prime, $a \in {\mathbb Z}$, then
$\left(\frac{a}{p}\right)$ equals $0$, if $p | a$; $1$ if $a$ is a quadratic
residue modulo $p$; and $-1$ otherwise.
Euler's criterion:
$\left(\frac{a}{p}\right)=a^{\left(\frac{p-1}{2}\right)} \pmod p$.

\subsubsection{Jacobi symbol}
If $n=p_1^{a_1} \cdots p_k^{a_k}$ is odd, then
$\left(\frac{a}{n}\right) = \prod_{i=1}^k \left(\frac{a}{p_i}\right)^{k_i}$.

\subsubsection{Primitive roots}
If the order of $g$ modulo $m$ (min $n>0$:
$g^n \equiv 1 \pmod{m}$) is $\phi(m)$, then $g$ is called a primitive root.
If $Z_m$ has a primitive root, then it has $\phi(\phi(m))$ distinct primitive
roots. $Z_m$ has a primitive root iff $m$ is one of $2$, $4$,
$p^k$, $2p^k$, where $p$ is an odd prime.
If $Z_m$ has a primitive root $g$, then for all $a$ coprime to $m$,
there exists unique integer $i=\text{ind}_g(a)$ modulo $\phi(m)$,
such that $g^i \equiv a \pmod{m}$.
$\text{ind}_g(a)$ has logarithm-like properties:
$\text{ind}(1) = 0$, $\text{ind}(ab) = \text{ind}(a) + \text{ind}(b)$.

If $p$ is prime and $a$ is not divisible by $p$, then congruence
$x^n \equiv a \pmod{p}$ has $\gcd(n, p-1)$ solutions if
$a^{(p-1)/\gcd(n,p-1)} \equiv 1 \pmod{p}$, and no solutions otherwise.

\subsubsection{Discrete logarithm problem}
Find $x$ from $a^x \equiv b \pmod{m}$.
Can be solved in $O(\sqrt{m})$ time and space with a meet-in-the-middle trick.
Let $n = \lceil \sqrt{m} \rceil$, and $x = ny - z$.
Equation becomes $a^{ny} \equiv b a^z \pmod{m}$.  Precompute all values that
the RHS can take for $z = 0, 1, \dots, n-1$, and brute force $y$ on the LHS,
each time checking whether there's a corresponding value for RHS.

\subsubsection{Pythagorean triples}
Integer solutions of $x^2 + y^2 = z^2$
All relatively prime triples are given by:
$x=2mn, y=m^2-n^2, z=m^2+n^2$ where $m>n, \gcd(m,n)=1$ and $m \not\equiv n \pmod{2}$.
All other triples are multiples of these.
Equation $x^2 + y^2 = 2z^2$ is equivalent to $(\frac{x+y}{2})^2 + (\frac{x-y}{2})^2 = z^2$.
The Pythagorean triples are uniquely generated by
\[ a=k\cdot (m^{2}-n^{2}),\ \,b=k\cdot (2mn),\ \,c=k\cdot (m^{2}+n^{2}), \]
with $m > n > 0$, $k > 0$, $m \bot n$, and either $m$ or $n$ even.

\subsubsection{Postage stamps/McNuggets problem}
Let $a$, $b$ be relatively-prime integers.
There are exactly $\frac{1}{2}(a-1)(b-1)$ numbers \emph{not} of form $ax+by$ ($x,y \ge 0$),
and the largest is $(a-1)(b-1)-1 = ab - a - b$.

\subsubsection{Fermat's two-squares theorem}
Odd prime $p$ can be represented
as a sum of two squares iff $p \equiv 1 \pmod 4$.
A product of two sums of two squares is a sum of two squares.
Thus, $n$ is a sum of two squares iff every prime of
form $p=4k+3$ occurs an even number of times in $n$'s factorization.

\subsubsection{Modular Arithmetic}
$i \bmod j = i - j\left\lfloor \frac{i}{j} \right\rfloor$
\begin{itemize}
  \item HCN: 1e6(240), 1e9(1344), 1e12(6720), 1e14(17280), 1e15(26880), 1e16(41472)
  \item \( \gcd(a, b, c, d, \dots) = \gcd(a, b - a, c - b, d - c, \dots) \)
  \item \( \gcd(a + k, b + k, c + k, d + k, \dots) = \gcd(a + k, b - a, c - b, d - c, \dots) \)
  \item Primitive root exists iff \( n = 1, 2, 4, p^k, 2\times p^k \), where \( p \) is an odd prime.
  \item If primitive root exists, there are \( \phi(\phi(n)) \) primitive roots of \( n \).
  \item The numbers from \( 1 \) to \( n \) have in total \( O(n\log\log n) \) unique prime factors.
  \item \( x \equiv r_1 \pmod{m_1} \) and \( x \equiv r_2 \pmod{m_2} \) has a solution iff \( \gcd(m_1, m_2) | (r_1 - r_2) \)
  \item \( ca \equiv cb \pmod{m} \iff a \equiv b \pmod{ \frac{m}{\gcd(m, c)}} \)
  \item \( ax \equiv b \pmod{m} \) has a solution \( \iff \) \( \gcd(a, m) | b \)
  \item If \( ax \equiv b \pmod{m} \) has a solution, then it has \( \gcd(a, m) \) solutions and they are separated by \( \frac{m}{\gcd(a, m)} \)
  \item \( ax \equiv 1 \pmod{m} \) has a solution or \( a \) is invertible \( \pmod{m} \) \( \iff \) \(\gcd(a, m) = 1 \)
  \item \( x^2 \equiv 1 \pmod{p} \) then \( x \equiv \pm 1 \pmod{p} \)
  \item When \( p \equiv 3 \pmod 4 \), solution of \( x^2 \equiv a \pmod p \) is \( x \equiv \pm a^{\frac{p + 1}{4}} \pmod p \)
  \item When \( p \equiv 5 \pmod 8 \), solution is \( x \equiv a^{\frac{p + 3}{8}} \pmod p \) or \( x \equiv 2^{\frac{p - 1}{4}} a^{\frac{p + 3}{8}} \pmod p \)
\end{itemize}

\subsubsection{Wilson's Theorem}
Wilson's Theorem states that an integer $p > 1$ is a prime number if and only if:
$$(p - 1)! \equiv -1 \pmod{p}$$

Equivalently, in positive modulo representation:
$$(p - 1)! \equiv p - 1 \pmod{p}$$

\textbf{Competitive Programming Applications:}
While not practical for general primality testing due to the $O(p)$ factorial computation, Wilson's Theorem is extremely useful for calculating factorials close to $p$ modulo $p$. 

If you need to find $(p-k)! \pmod{p}$ for a small $k$, you can use Wilson's Theorem to compute it in $O(k \log p)$ time (using modular inverse) instead of $O(p)$ time.

\textbf{Useful Corollaries:}
By expanding $(p-1)! = (p-1)(p-2)!$ and dividing out terms, we get:
\begin{itemize}
  \item $(p - 1)! \equiv -1 \pmod{p}$
  \item $(p - 2)! \equiv 1 \pmod{p}$
  \item $(p - k)! \equiv \frac{(-1)^k}{(k-1)!} \pmod{p}$ \quad (for $1 \le k \le p$)
\end{itemize}

\subsection{Combinatorics}
\subsubsection{Permutations and Combinations Formulas}
Permutation: ${}^n P_r = \frac{n!}{(n-r)!}$ \\
Combination: ${}^n C_r = \frac{n!}{r!(n-r)!}$

\subsubsection{Factorial}
\begin{center}
\begin{tabular}{l}
  \begin{tabular}{c|c@{\ }c@{\ }c@{\ }c@{\ }c@{\ }c@{\ }c@{\ }c@{\ }c@{\ }c}
    $n$  & 1 & 2 & 3 & 4  & 5   & 6   & 7    & 8     & 9      & 10\\
    \hline
    $n!$ & 1 & 2 & 6 & 24 & 120 & 720 & 5040 & 40320 & 362880 & 3628800\\
  \end{tabular}\\
  \begin{tabular}{c|c@{\ }c@{\ }c@{\ }c@{\ }c@{\ }c@{\ }c@{\ }c@{\ }c@{\ }c}
    $n$  & 11    & 12    & 13    & 14     & 15     & 16     & 17\\
    \hline
    $n!$ & 4.0e7 & 4.8e8 & 6.2e9 & 8.7e10 & 1.3e12 & 2.1e13 & 3.6e14\\
  \end{tabular}\\
  \begin{tabular}{c|c@{\ }c@{\ }c@{\ }c@{\ }c@{\ }c@{\ }c@{\ }c@{\ }c@{\ }c}
    $n$  & 20   & 25   & 30   & 40   & 50   & 100   & 150   & 171\\
    \hline
    $n!$ & 2e18 & 2e25 & 3e32 & 8e47 & 3e64 & 9e157 & 6e262 & \scriptsize{$>$DBL\_MAX}\\
  \end{tabular}
\end{tabular}
\end{center}

\subsubsection{Binomial Coefficient}
\begin{itemize}
  \item Factoring in: \( \binom{n}{k} = \frac{n}{k} \binom{n - 1}{k - 1} \)
  \item Sum over \( k \): \( \sum_{k = 0}^n \binom{n}{k} = 2^n \)
  \item Alternating sum: \( \sum_{k = 0}^n (-1)^k \binom{n}{k} = 0 \)
  \item Even and odd sum: \( \sum_{k = 0}^n \binom{n}{2k} = \sum_{k = 0}^n \binom{n}{2k + 1} = 2^{n - 1} \)
  \item The Hockey Stick Identity:
    \begin{itemize}
      \item (Left to right) Sum over \( n \) and \( k \): \( \sum_{k = 0}^m \binom{n + k}{k} = \binom{n + m + 1}{m} \)
      \item (Right to left) Sum over \( n \): \( \sum_{m = k}^n \binom{m}{k} = \binom{n + 1}{k + 1} \)
    \end{itemize}
  \item Sum of the squares: \( \sum_{k = 0}^n \binom{n}{k}^2 = \binom{2n}{n} \)
  \item Weighted sum: \( \sum_{k = 1}^n k \binom{n}{k} = n2^{n - 1} \)
  \item Connection with the Fibonacci numbers: \( \sum_{k = 0}^{\lfloor n/2 \rfloor}\binom{n - k}{k} = F_{n + 1} \)
  \item Vandermonde's Identity: \( \sum_{i = 0}^k \binom{m}{i} \binom{n}{k - i} = \binom{m + n}{k} \)
  \item If \( f(n, k) = \binom{n}{0} + \binom{n}{1} + \dots + \binom{n}{k} \), Then \( f(n + 1, k) = 2 f(n, k) - \binom{n}{k} \) [For multiple \( f(n, k) \) queries, use Mo's algo]
\end{itemize}

\textbf{Lucas Theorem}
\[ \binom{m}{n} \equiv \prod_{i = 0}^k \binom{m_i}{n_i} \pmod p \]
\begin{itemize}
  \item \( \binom{m}{n} \) is divisible by $p$ if and only if at least one of the base-\(p\) digits of \( n \) is greater than the corresponding base-\( p \) digit of \( m \).
  \item The number of entries in the \( n \)-th row of Pascal's triangle that are not divisible by \( p = \prod_{i = 0}^k (n_i + 1) \)
  \item All entries in the \( (p^{k} - 1) \)-th row are not divisible by \( p \).
  \item \( \binom{n}{m} \equiv \lfloor \frac{n}{p} \rfloor \pmod p \)
\end{itemize}

\subsubsection{Cycles}
Let $g_S(n)$ be the number of $n$-permutations whose cycle lengths all belong to the set $S$. Then
$$\sum_{n=0} ^\infty g_S(n) \frac{x^n}{n!} = \exp\left(\sum_{n\in S} \frac{x^n} {n} \right)$$

\subsubsection{Derangements}
Permutations of a set such that none of the elements appear in their original position.
\[ \mkern-2mu D(n) = (n-1)(D(n-1)+D(n-2)) = n D(n-1)+(-1)^n = \left\lfloor\frac{n!}{e}\right\rceil \]

\subsubsection{Burnside's lemma}
Given a group $G$ of symmetries and a set $X$, the number of elements of $X$ \emph{up to symmetry} equals
\[ {\frac {1}{|G|}}\sum _{{g\in G}}|X^{g}|, \]
where $X^{g}$ are the elements fixed by $g$ ($g.x = x$).

If $f(n)$ counts ``configurations'' (of some sort) of length $n$, we can ignore rotational symmetry using $G = \mathbb Z_n$ to get
\[ g(n) = \frac 1 n \sum_{k=0}^{n-1}{f(\text{gcd}(n, k))} = \frac 1 n \sum_{k|n}{f(k)\phi(n/k)} \]

\subsubsection{Multinomial Coefficients}
\[
\frac{(C_1 + C_2 + \cdots + C_N)!}{C_1! \, C_2! \cdots C_N!}
=
\binom{C_1}{C_1}
\binom{C_1 + C_2}{C_2}
\binom{C_1 + C_2 + C_3}{C_3}
\cdots
\binom{C_1 + \cdots + C_N}{C_N}
\]

\subsubsection{Partitions and Subsets}
\textbf{Partition function}
Number of ways of writing $n$ as a sum of positive integers, disregarding the order of the summands.
\[ p(0) = 1,\ p(n) = \sum_{k \in \mathbb Z \setminus \{0\}}{(-1)^{k+1} p(n - k(3k-1) / 2)} \]
\[ p(n) \sim \frac{0.145}{n} \exp(2.56 \sqrt{n}) \]

\begin{center}
\begin{tabular}{c|c@{\ }c@{\ }c@{\ }c@{\ }c@{\ }c@{\ }c@{\ }c@{\ }c@{\ }c@{\ }c@{\ }c@{\ }c}
  $n$    & 0 & 1 & 2 & 3 & 4 & 5 & 6  & 7  & 8  & 9  & 20  & 50  & 100 \\ \hline
  $p(n)$ & 1 & 1 & 2 & 3 & 5 & 7 & 11 & 15 & 22 & 30 & 627 & $\mathtt{\sim}$2e5 & $\mathtt{\sim}$2e8 \\
\end{tabular}
\end{center}

\textbf{Partition Number Calculation}
\begin{itemize}
\item Time Complexity: $ O(n\sqrt{n}) $
\end{itemize}
\begin{verbatim}
for (int i = 1; i <= n; ++i) {
  pent[2 * i - 1] = i * (3 * i - 1) / 2;
  pent[2 * i] = i * (3 * i + 1) / 2;
}
p[0] = 1;
for (int i = 1; i <= n; ++i) {
  p[i] = 0;
  for (int j = 1, k = 0; pent[j] <= i; ++j) {
    if (k < 2) p[i] = add(p[i], p[i - pent[j]]);
    else p[i] = sub(p[i], p[i - pent[j]]); ++k, k &= 3;
  }
}
\end{verbatim}
The number of partitions of a positive integer \( n \) into exactly \( k \) parts equals the number of partitions of \( n \) whose largest part equals \( k \)
\[ p_k(n) = p_k(n - k) + p_{k - 1}(n - 1) \]

\textbf{2nd Kaplansky's Lemma}
The number of ways of selecting \( k \) objects, no two consecutive, from \( n \) labelled objects arrayed in a circle is \( \frac{n}{k} \binom{n-k-1}{k-1} = \frac{n}{n - k} \binom{n-k}{k} \)

\textbf{Distinct Objects into Distinct Bins}
\begin{itemize}
  \item $ n $ distinct objects into $ r $ distinct bins $ = r^n $
  \item Among $ n $ distinct objects, exactly $ k $ of them into $ r $ distinct bins $ = \binom{n}{k}r^k $
  \item $ n $ distinct objects into $ r $ distinct bins such that each bin contains at least one object $ = \sum_{i = 0}^{r} (-1)^i \binom{r}{i} (r - i)^n $
\end{itemize}

\subsubsection{General purpose numbers}
\textbf{Eulerian numbers}
Number of permutations $\pi \in S_n$ in which exactly $k$ elements are greater than the previous element.
$$E(n,k) = (n-k)E(n-1,k-1) + (k+1)E(n-1,k)$$
$$E(n,0) = E(n,n-1) = 1$$
$$E(n,k) = \sum_{j=0}^k(-1)^j\binom{n+1}{j}(k+1-j)^n$$

\textbf{Bell numbers}
Total number of partitions of $n$ distinct elements. $B(n) = 1, 1, 2, 5, 15, 52, 203, 877, 4140, 21147, \dots$. For $p$ prime,
\[ B(p^m+n)\equiv mB(n)+B(n+1) \pmod{p} \]

\textbf{Bernoulli numbers}
$\Scale[.95]{\sum\limits_{j=0}^m \binom{m+1}{j} B_j = 0}$.
\quad $B_0=1$, $B_1=-\frac{1}{2}$. $B_n=0$, for all odd $n \ne 1$.

\textbf{Catalan numbers}
\[ C_n=\frac{1}{n+1}\binom{2n}{n}= \binom{2n}{n}-\binom{2n}{n+1} = \frac{(2n)!}{(n+1)!n!} \]
\[ C_0=1,\ C_{n+1} = \frac{2(2n+1)}{n+2}C_n,\ C_{n+1}=\sum C_iC_{n-i} \]
\begin{itemize}[noitemsep]
\item $C_n = 1, 1, 2, 5, 14, 42, 132, 429, 1430, 4862, 16796, 58786, \dots$
\item sub-diagonal monotone paths in an $n\times n$ grid.
\item strings with $n$ pairs of parenthesis, correctly nested.
\item binary trees with $n+1$ leaves (0 or 2 children).
\item ordered trees with $n+1$ vertices.
\item ways a convex polygon with $n+2$ sides can be cut into triangles by connecting vertices with straight lines.
\item permutations of $[n]$ with no 3-term increasing subseq.
\item Find the count of balanced parentheses sequences consisting of $n + k$ pairs of parentheses where the first $k$ symbols are open brackets.
  $$C_n^{(k)} = \frac{k + 1}{n + k + 1} \binom{2n + k}{n}$$
\item Recursive formula of Catalan Numbers:
  $$ C_{n}^{(k)} = \frac{(2n + k - 1) \cdot (2n + k)}{n \cdot (n + k + 1)} C_{n - 1}^{(k)} $$
\end{itemize}

\textbf{Lucas Number}
Number of edge cover of a cycle graph $ C_n $ is $ L_n $
$$ L(n) = L(n-1) + L(n-2); L(0)=2, L(1)=1 $$

\subsubsection{Stars and Bars Theorems}
\begin{enumerate}
  \item No Empty Bins Allowed: $\dbinom{N - 1}{K - 1}$
  \item Empty Bins Allowed: $\dbinom{N + K - 1}{K - 1}$
\end{enumerate}

\subsubsection{Ballot Theorem}
Suppose that in an election, candidate A receives $a$ votes and candidate B receives $b$ votes, where $a \ge kb$ for some positive integer $k$. Compute the number of ways the ballots can be ordered so that A maintains more than $k$ times as many votes as B throughout the counting of the ballots.
The solution to the ballot problem is $\frac{a - kb}{a+b} \times \binom{a + b}{a}$

\subsection{Sequences and Series}
\subsubsection{Fibonacci Number}
\begin{enumerate}
  \item $k=A-B, F_A F_B = F_{k+1} F_A^2 + F_k F_A F_{A-1}$
  \item $\sum_{i=0}^n F_i^2 = F_{n+1} F_n$
  \item $\sum_{i=0}^n F_i F_{i+1} = F_{n+1}^2 - (-1)^n$
  \item $\sum_{i=0}^n F_i F_{i-1} = \sum_{i=0}^{n-1} F_i F_{i+1}$
  \item $\operatorname{gcd}\left(F_m, F_n\right) = F_{\operatorname{gcd}(m, n)}$
  \item $\operatorname{gcd}\left(F_n, F_{n+1}\right) = \operatorname{gcd}\left(F_n, F_{n+2}\right) = \operatorname{gcd}\left(F_{n+1}, F_{n+2}\right) = 1$
\end{enumerate}

\subsubsection{Sums}
\begin{align*}
1 + 2 + 3 + \dots + n &= \frac{n(n+1)}{2} \\
1 + 3 + 5 + \dots + (2n-1) &= n^2 \\
2 + 4 + 6 + \dots + 2n &= n(n+1) \\
1^2 + 2^2 + 3^2 + \dots + n^2 &= \frac{n(n+1)(2n+1)}{6} \\
1^3 + 2^3 + 3^3 + \dots + n^3 &= \frac{n^2(n+1)^2}{4} \\
1^4 + 2^4 + 3^4 + \dots + n^4 &= \frac{n(n+1)(2n+1)(3n^2 + 3n - 1)}{30}
\end{align*}

$\Scale[0.85]{ \sum\limits_{i=1}^{n} i^{m} = \frac{1}{m + 1} \left[ (n + 1)^{m + 1} - 1 - \sum\limits_{i=1}^{n} \left((i+1)^{m+1} - i^{m+1} - (m + 1)i^{m}\right)\right] }$ \\
$\Scale[0.99]{\sum\limits_{i=1}^{n-1} i^{m} = \frac{1}{m + 1} \sum\limits_{k=0}^{m} \binom{m+1}{k} B_{k}n^{m + 1 - k}}$ \\
$\sum\limits_{k=0}^n kx^k = \frac{x - (n+1)x^{n+1} + nx^{n+2}}{(x-1)^2}$

\subsubsection{Progressions}
\textbf{Arithmetic Progression}
\begin{itemize}
\item \textbf{Sequence:} $a, a+d, a+2d, \dots, a+(n-1)d$
\item \textbf{Sum of first $n$ terms:} $S_n = \frac{n}{2} [2a + (n-1)d]$
\end{itemize}

\textbf{Geometric Progression}
\begin{itemize}
  \item \textbf{Sequence:} $a, ar, ar^2, \dots, ar^{n-1}$
  \item \textbf{Sum (for $r > 1$):} $S_n = \frac{a(r^n - 1)}{r-1}$
  \item \textbf{Sum (for $r < 1$):} $S_n = \frac{a(1 - r^n)}{1-r}$
\end{itemize}

\subsubsection{Series}
$$e^x = 1+x+\frac{x^2}{2!}+\frac{x^3}{3!}+\dots,\,(-\infty<x<\infty)$$
$$\ln(1+x) = x-\frac{x^2}{2}+\frac{x^3}{3}-\frac{x^4}{4}+\dots,\,(-1<x\le1)$$
$$\sqrt{1+x} = 1+\frac{x}{2}-\frac{x^2}{8}+\frac{2x^3}{32}-\frac{5x^4}{128}+\dots,\,(-1\le x\le1)$$
$$\sin x = x-\frac{x^3}{3!}+\frac{x^5}{5!}-\frac{x^7}{7!}+\dots,\,(-\infty<x<\infty)$$
$$\cos x = 1-\frac{x^2}{2!}+\frac{x^4}{4!}-\frac{x^6}{6!}+\dots,\,(-\infty<x<\infty)$$
$$(x + a)^{-n} = \sum\limits_{k=0}^{\infty} (-1)^{k} \binom{n + k - 1}{k} x^{k}a^{-n-k}$$

\textbf{Generating Function}
$$ 1/(1 - x) = 1 + x + x^2 + x^3 + \dots $$
$$ 1/(1 - ax) = 1 + ax + (ax)^2 + (ax)^3 + \dots $$
$$ 1/(1 - x)^2 = 1 + 2x + 3x^2 + 4x^3 + \dots $$
$$ 1/(1 - x)^3 = \binom{2}{2} + \binom{3}{2}x + \binom{4}{2}x^2 + \binom{5}{2}x^3 + \dots $$
$$ 1/(1 - ax)^{k + 1} = 1 + \binom{1 + k}{k}(ax) + \binom{2 + k}{k}(ax)^2 + \binom{3 + k}{k}(ax)^3 + \dots $$
$$ x(x + 1)(1 - x)^{-3} = 1 + x + 4x^2 + 9x^3 + 16x^4 + 25x^5 + \dots $$
$$ e^x = 1 + x + x^2/2! + x^3/3! + x^4/4! + \dots $$

\subsubsection{Logarithms}
\textbf{Change of Base Formula}
$$
\log_a x = \frac{\log_b x}{\log_b a}
$$

\subsubsection{Inequalities}
\textbf{Titu's Lemma}
For positive reals $a_1, a_2, \ldots, a_n$ and $b_1, b_2, \ldots, b_n$,
$$\frac{{a_1}^2}{b_1} + \frac{{a_2}^2}{b_2} + \dots + \frac{{a_n}^2}{b_n} \ge \frac{(a_1 + a_2 + \dots + a_n)^2}{b_1 + b_2 + \dots + b_n}$$
Equality holds if and only if $ a_i = kb_i$ for a non-zero real constant $k$.

\subsection{Graph Theory}

\subsubsection{Coloring}
\begin{itemize}
\item The number of labeled undirected graphs with $ n $ vertices, $ G_n = 2^{\binom{n}{2}} $
\item The number of labeled directed graphs with $ n $ vertices, $ G_n = 2^{n(n-1)} $
\item The number of connected labeled undirected graphs with $ n $ vertices, $ C_n = 2^{\binom{n}{2}} - \frac{1}{n} \sum_{k = 1}^{n - 1} k \binom{n}{k} 2^{\binom{n-k}{2}}C_k = 2^{\binom{n}{2}} - \sum_{k = 1}^{n - 1} \binom{n - 1}{k - 1} 2^{\binom{n-k}{2}}C_k $
\item The number of k-connected labeled undirected graphs with $ n $ vertices, $ D[n][k] = \sum_{s = 1}^{n} \binom{n - 1}{s- 1}C_s D[n - s][k - 1] $
\item Cayley's formula: the number of trees on $ n $ labeled vertices = the number of spanning trees of a complete graph with $ n $ labeled vertices = $ n^{n - 2} $
\item Number of ways to color a graph using $ k $ colors such that no two adjacent nodes have the same color:
  \begin{itemize}
    \item Complete graph = $ k(k-1)(k-2)\dots(k-n+1) $
    \item Tree = $ k(k - 1)^{n - 1} $
    \item Cycle = $ (k - 1)^n + (-1)^n (k - 1) $
  \end{itemize}
\item Number of trees with $ n $ labeled nodes: $ n^{n - 2} $
\end{itemize}

\subsubsection{Classical Problem}
$ F(n, k) = $ number of ways to color $n$ objects using exactly $ k $ colors

Let $ G(n, k) $ be the number of ways to color $n$ objects using no more than $ k $ colors.

Then, $ F(n, k) = G(n, k) - \binom{k}{1} G(n, k-1) + \binom{k}{2} G(n, k-2) - \binom{k}{3} G(n, k-3) \dots $

\textbf{Determining G(n, k)} :
Suppose, we are given a $1 \times n$ grid. Any two adjacent cells can not have the same color.
Then, $G(n, k) = k (k-1)^{n-1}$

If no such condition on adjacent cells:
Then, $G(n, k) = k^n$

\subsubsection{Matching Formula}
\textbf{Normal Graph} \\
MM + MEC = n (excluding vertex), IS + VC = G, MIS + MVC = G

\textbf{Bipartite Graph} \\
MIS = n - MBM, MVC = MBM, MEC = n - MBM

\subsection{Strings and Bitwise}
\subsubsection{String}
\begin{itemize}
\item If the sum of length of some strings is $ N $, there can be at most $ \sqrt{N} $ distinct lengths.
\item A Text can have at most $ O(N \sqrt{N}) $ distinct substrings that match with given patterns where the sum of the length of the given patterns is $ N $.
\item Period =  $n \pmod{n - \text{pi.back()}} == 0 \text{ ? } n - \text{pi.back()} : n$
\item The first (period) cyclic rotations of a string are distinct. Further cyclic rotations repeat the previous strings.
\item $ S $ is a palindrome if and only if its period is a palindrome.
\item If $S$ and $T$ are palindromes, then the periods of $S$ \& $T$ are the same if and only if $S + T$ is a palindrome.
\end{itemize}

\subsubsection{Tree Hashing}
$ f(u) = sz[u] \times \sum_{i = 0} f(v) \times p^{i} $;    $ f(v) $ are sorted \\
$ f(child) = 1 $

\subsubsection{Bit}
\begin{itemize}
\item $(a \oplus b)$ and $(a + b)$ have the same parity
\item $(a + b) = (a \oplus b) + 2(a \ \& \ b)$
\item $\gcd(a, b) \le |a - b| \le (a \oplus b)$
\end{itemize}

\subsubsection{Convolution}
\begin{itemize}
\item Hamming Distance: Replace $ 0 $ with $ -1 $
\item SQRT Decomposition: Find block size, $B = \sqrt{8n}$
\end{itemize}

\subsubsection{Permutation Adjacency Difference}
To maximize the sum of adjacent differences of a permutation, it is necessary and sufficient to place the smallest half numbers in odd positions and the greatest half numbers in even positions. Or, vice versa.

\subsection{Game Theory}
\subsubsection{Grundy numbers}
For a two-player, normal-play (last to move wins) game on a graph $(V,E)$:
$G(x) = \mbox{mex}(\{ G(y) : (x, y) \in E \})$,
where $\mbox{mex}(S) = \min \{ n \ge 0: n \not\in S \}$.
$x$ is losing iff $G(x) = 0$.

\subsubsection{Sums of games}
\begin{itemize}
\item
  \emph{Player chooses a game and makes a move in it:}
  Grundy number of a position is xor of grundy numbers of positions in summed games.
\item
  \emph{Player chooses a non-empty subset of games (possibly, all) and makes moves in all of them:}
  A position is losing iff each game is in a losing position.
\item
  \emph{Player chooses a proper subset of games (not empty and not all), and makes moves in all chosen ones:}
  A position is losing iff grundy numbers of all games are equal.
\item
  \emph{Player must move in all games, and loses if can't move in some game:}
  A position is losing if any of the games is in a losing position.
\end{itemize}

\subsubsection{Mis\`{e}re Nim}
A position with pile sizes $a_1, a_2, \dots, a_n \ge 1$,
not all equal to $1$, is losing iff $a_1 \oplus a_2 \dots \oplus a_n = 0$
(like in normal nim.)
A position with $n$ piles of size $1$ is losing iff $n$ is \emph{odd}.

\end{document}

\end{multicols}
\end{document}